\documentclass{article}
\usepackage{amsmath,amssymb,amsthm}
\IfFileExists{hyperref.sty}{%
  \IfFileExists{color.sty}{\usepackage{color}\definecolor{paperlinkblue}{rgb}{0.10,0.25,0.42}%
    \usepackage[colorlinks=true,linkcolor=paperlinkblue]{hyperref}}{\usepackage[hidelinks]{hyperref}}%
}{}
\ifdefined\hyperlink
  \newcommand{\paperanchor}[1]{\hypertarget{#1}{}}%
  \newcommand{\paperlink}[2]{\hyperlink{#1}{\begin{NoHyper}#2\end{NoHyper}}}%
\else
  \ifdefined\pdfextension
    \ifnum\outputmode>0
      \def\paperdest{\pdfextension dest}\def\paperstart{\pdfextension startlink}%
      \def\paperend{\pdfextension endlink}\def\papercolor{\pdfextension colorstack}%
    \fi
  \else\ifdefined\pdfstartlink
    \ifnum\pdfoutput>0
      \let\paperdest\pdfdest\let\paperstart\pdfstartlink\let\paperend\pdfendlink\let\papercolor\pdfcolorstack
    \fi
  \fi\fi
  \ifdefined\paperdest
    \newcommand{\paperanchor}[1]{\paperdest name{#1} xyz}%
    \newcommand{\paperlink}[2]{\leavevmode\paperstart attr{/Border [0 0 0]} goto name{#1}%
      \papercolor 0 push {0.10 0.25 0.42 rg 0.10 0.25 0.42 RG}#2\papercolor 0 pop\paperend}%
  \else\ifdefined\XeTeXversion
    \newcommand{\paperanchor}[1]{\special{pdf:dest (#1) [@thispage /XYZ null @ypos null]}}%
    \newcommand{\paperlink}[2]{\leavevmode\special{pdf:bann << /Subtype /Link /Border [0 0 0] /A << /S /GoTo /D (#1) >> >>}%
      \special{color push rgb 0.10 0.25 0.42}#2\special{color pop}\special{pdf:eann}}%
  \else
    \PackageError{lean-to-latex}{No PDF hyperlink backend available}{Use XeLaTeX, pdfLaTeX, LuaLaTeX, or install hyperref.}%
  \fi\fi
\fi
\newcommand{\paperlineanchor}[1]{\raisebox{10pt}[0pt][0pt]{\paperanchor{#1}}}
\DeclareRobustCommand{\paperresultname}[1]{\paperlineanchor{theorem:\arabic{theorem}}#1}
\newtheoremstyle{paperresult}{8pt}{5pt}{\normalfont\raggedright}{}{\bfseries}{.}{\newline}{}
\theoremstyle{paperresult}
\newtheorem{theorem}{\paperresultname{Theorem}}
\newcounter{paperpremise}
\newtheorem{lemma}[theorem]{\paperresultname{Lemma}}
\newtheorem{corollary}[theorem]{\paperresultname{Corollary}}
\newcounter{definingequation}
\newenvironment{proofstep}{\par\smallskip\ifdim\leftskip<3em\relax\advance\leftskip by 1em\relax\fi\noindent\ignorespaces}{\par\smallskip}
\setlength{\textwidth}{6.5in}
\setlength{\oddsidemargin}{0in}
\AddToHook{env/proof/begin}{\raggedright}
\setlength{\emergencystretch}{\linewidth}
\title{Intersecting families of sets}
\author{}
\date{}
\begin{document}
\maketitle
\section*{Definitions}
\paragraph{Intersecting.} \par\noindent\refstepcounter{definingequation}\label{defeq:1}\paperlineanchor{defeq:1}\textup{(\paperlink{defeq:1}{D1})} Let \(\alpha\) be a type. Fix a structure \(\mathit{inst}\) of type \(\operatorname{Semilattice\,Inf}\left(\alpha\right)\). Fix a structure \(\mathit{inst}_{1}\) of type \(\operatorname{Order\,Bot}\left(\alpha\right)\). Let \(s\) be a subset of \(\alpha\). Then \(\operatorname{Intersecting}(s) = (\forall a\, :\, \alpha,\; a \in s \implies (\forall b\, :\, \alpha,\; b \in s \implies \neg \operatorname{Disjoint}(a,\allowbreak{} b)))\).

\subsection*{Prerequisites}
\begin{description}
\item[\paperlineanchor{premise:416e642e696d705f6c656674}\paperlink{premise:416e642e696d705f6c656674}{P1.}] \refstepcounter{paperpremise}\label{premise:416e642e696d705f6c656674}\textbf{Imp left.} Let \(a\), \(b\), and \(c\) be propositions. Assume \(a \implies b\). Assume \(a \land c\). Then \(b \land c\).
\item[\paperlineanchor{premise:436c6173736963616c2e6e6f745f6e6f74}\paperlink{premise:436c6173736963616c2e6e6f745f6e6f74}{P2.}] \refstepcounter{paperpremise}\label{premise:436c6173736963616c2e6e6f745f6e6f74}\textbf{Not not.} Let \(a\) be a proposition. Then \(\neg \left(\neg a\right) \iff a\).
\item[\paperlineanchor{premise:4469736a6f696e742e6c655f636f6d706c5f6c656674}\paperlink{premise:4469736a6f696e742e6c655f636f6d706c5f6c656674}{P3.}] \refstepcounter{paperpremise}\label{premise:4469736a6f696e742e6c655f636f6d706c5f6c656674}\textbf{At most compl left.} Let \(\alpha\) be a type. Fix a structure \(\mathit{inst}\) of type \(\operatorname{Heyting\,Algebra}\left(\alpha\right)\). Let \(a\) and \(b\) be elements of \(\alpha\). Assume \(\operatorname{Disjoint}\left(b, a\right)\). Then \(a \le \operatorname{compl}\left(b\right)\).
\item[\paperlineanchor{premise:4469736a6f696e742e6d6f6e6f5f6c656674}\paperlink{premise:4469736a6f696e742e6d6f6e6f5f6c656674}{P4.}] \refstepcounter{paperpremise}\label{premise:4469736a6f696e742e6d6f6e6f5f6c656674}\textbf{Mono left.} Let \(\alpha\) be a type. Equip \(\alpha\) with a partial order. Fix a structure \(\mathit{inst}_{1}\) of type \(\operatorname{Order\,Bot}\left(\alpha\right)\). Let \(a\), \(b\), and \(c\) be elements of \(\alpha\). Assume \(a \le b\). Assume \(\operatorname{Disjoint}\left(b, c\right)\). Then \(\operatorname{Disjoint}\left(a, c\right)\).
\item[\paperlineanchor{premise:4469736a6f696e742e73796d6d}\paperlink{premise:4469736a6f696e742e73796d6d}{P5.}] \refstepcounter{paperpremise}\label{premise:4469736a6f696e742e73796d6d}\textbf{Symm.} Let \(\alpha\) be a type. Equip \(\alpha\) with a partial order. Fix a structure \(\mathit{inst}_{1}\) of type \(\operatorname{Order\,Bot}\left(\alpha\right)\). Let \(a\) and \(b\) be elements of \(\alpha\). Assume \(\operatorname{Disjoint}\left(a, b\right)\). Then \(\operatorname{Disjoint}\left(b, a\right)\).
\item[\paperlineanchor{premise:4578697374732e696d70}\paperlink{premise:4578697374732e696d70}{P6.}] \refstepcounter{paperpremise}\label{premise:4578697374732e696d70}\textbf{Imp.} Let \(\alpha\) be a sort. Let \(p\) and \(q\) be functions of type \(\alpha \to \mathsf{Prop}\). Assume \(\forall a\, :\, \alpha,\; p\left(a\right) \implies q\left(a\right)\). Assume \(\exists a_{1}\, :\, \alpha,\; p\left(a_{1}\right)\). Then \(\exists a_{1}\, :\, \alpha,\; q\left(a_{1}\right)\).
\item[\paperlineanchor{premise:46696e7365742e636172645f6469736a556e696f6e}\paperlink{premise:46696e7365742e636172645f6469736a556e696f6e}{P7.}] \refstepcounter{paperpremise}\label{premise:46696e7365742e636172645f6469736a556e696f6e}\textbf{Cardinality disj union.} Let \(\alpha\) be a type. Let \(s\) and \(t\) be finite subsets of \(\alpha\). Let \(h\) be a proof of \(\operatorname{Disjoint}\left(s, t\right)\). Then \(\left|\operatorname{disj\,Union}\left(s, t\right)\right| = \left|s\right| + \left|t\right|\).
\item[\paperlineanchor{premise:46696e7365742e636172645f6c655f756e6976}\paperlink{premise:46696e7365742e636172645f6c655f756e6976}{P8.}] \refstepcounter{paperpremise}\label{premise:46696e7365742e636172645f6c655f756e6976}\textbf{Cardinality at most whole space.} Let \(\alpha\) be a type. Fix a finite enumeration of \(\alpha\). Let \(s\) be a finite subset of \(\alpha\). Then \(\left|s\right| \le \left|\alpha\right|\).
\item[\paperlineanchor{premise:46696e7365742e636172645f6d6170}\paperlink{premise:46696e7365742e636172645f6d6170}{P9.}] \refstepcounter{paperpremise}\label{premise:46696e7365742e636172645f6d6170}\textbf{Cardinality map.} Let \(\alpha\) and \(\beta\) be types. Let \(s\) be a finite subset of \(\alpha\). Let \(f\) be an element of \(\operatorname{Embedding}\left(\alpha, \beta\right)\). Then \(\left|\operatorname{map}\left(f, s\right)\right| = \left|s\right|\).
\item[\paperlineanchor{premise:46696e7365742e636f655f65715f656d707479}\paperlink{premise:46696e7365742e636f655f65715f656d707479}{P10.}] \refstepcounter{paperpremise}\label{premise:46696e7365742e636f655f65715f656d707479}\textbf{Coercion equals empty.} Let \(\alpha\) be a type. Let \(s\) be a finite subset of \(\alpha\). Then \(\operatorname{underlying\,set}(s) = \varnothing \iff s = \varnothing\).
\item[\paperlineanchor{premise:46696e7365742e636f655f65715f756e6976}\paperlink{premise:46696e7365742e636f655f65715f756e6976}{P11.}] \refstepcounter{paperpremise}\label{premise:46696e7365742e636f655f65715f756e6976}\textbf{Coercion equals whole space.} Let \(\alpha\) be a type. Fix a finite enumeration of \(\alpha\). Let \(s\) be a finite subset of \(\alpha\). Then \(\operatorname{underlying\,set}(s) = \alpha \iff s = \alpha\).
\item[\paperlineanchor{premise:46696e7365742e636f655f696e73657274}\paperlink{premise:46696e7365742e636f655f696e73657274}{P12.}] \refstepcounter{paperpremise}\label{premise:46696e7365742e636f655f696e73657274}\textbf{Coercion insert.} Let \(\alpha\) be a type. Assume equality on \(\alpha\) is decidable. Let \(a\) be an element of \(\alpha\). Let \(s\) be a finite subset of \(\alpha\). Then \(\operatorname{underlying\,set}(\{a\} \cup s) = \{a\} \cup \operatorname{underlying\,set}(s)\).
\item[\paperlineanchor{premise:46696e7365742e636f655f6d6170}\paperlink{premise:46696e7365742e636f655f6d6170}{P13.}] \refstepcounter{paperpremise}\label{premise:46696e7365742e636f655f6d6170}\textbf{Coercion map.} Let \(\alpha\) and \(\beta\) be types. Let \(f\) be an element of \(\operatorname{Embedding}\left(\alpha, \beta\right)\). Let \(s\) be a finite subset of \(\alpha\). Then \(\operatorname{underlying\,set}(\operatorname{map}(f,\allowbreak{} s)) = \operatorname{coe}(f)[\operatorname{underlying\,set}(s)]\).
\item[\paperlineanchor{premise:46696e7365742e636f655f73696e676c65746f6e}\paperlink{premise:46696e7365742e636f655f73696e676c65746f6e}{P14.}] \refstepcounter{paperpremise}\label{premise:46696e7365742e636f655f73696e676c65746f6e}\textbf{Coercion singleton.} Let \(\alpha\) be a type. Let \(a\) be an element of \(\alpha\). Then \(\operatorname{underlying\,set}(\{a\}) = \{a\}\).
\item[\paperlineanchor{premise:46696e7365742e636f655f756e696f6e}\paperlink{premise:46696e7365742e636f655f756e696f6e}{P15.}] \refstepcounter{paperpremise}\label{premise:46696e7365742e636f655f756e696f6e}\textbf{Coercion union.} Let \(\alpha\) be a type. Assume equality on \(\alpha\) is decidable. Let \(s_{1}\) and \(s_{2}\) be finite subsets of \(\alpha\). Then \(\operatorname{underlying\,set}(s_{1} \cup s_{2}) = \operatorname{underlying\,set}(s_{1}) \cup \operatorname{underlying\,set}(s_{2})\).
\item[\paperlineanchor{premise:46696e7365742e6469736a556e696f6e5f65715f756e696f6e}\paperlink{premise:46696e7365742e6469736a556e696f6e5f65715f756e696f6e}{P16.}] \refstepcounter{paperpremise}\label{premise:46696e7365742e6469736a556e696f6e5f65715f756e696f6e}\textbf{Disj union equals union.} Let \(\alpha\) be a type. Assume equality on \(\alpha\) is decidable. Let \(s\) and \(t\) be finite subsets of \(\alpha\). Let \(h\) be a proof of \(\operatorname{Disjoint}\left(s, t\right)\). Then \(\operatorname{disj\,Union}\left(s, t\right) = s \cup t\).
\item[\paperlineanchor{premise:46696e7365742e6469736a6f696e745f636f65}\paperlink{premise:46696e7365742e6469736a6f696e745f636f65}{P17.}] \refstepcounter{paperpremise}\label{premise:46696e7365742e6469736a6f696e745f636f65}\textbf{Disjoint coercion.} Let \(\alpha\) be a type. Let \(s\) and \(t\) be finite subsets of \(\alpha\). Then \(\operatorname{Disjoint}(\operatorname{underlying\,set}(s),\allowbreak{} \operatorname{underlying\,set}(t)) \iff \operatorname{Disjoint}(s,\allowbreak{} t)\).
\item[\paperlineanchor{premise:46696e7365742e6469736a6f696e745f6c656674}\paperlink{premise:46696e7365742e6469736a6f696e745f6c656674}{P18.}] \refstepcounter{paperpremise}\label{premise:46696e7365742e6469736a6f696e745f6c656674}\textbf{Disjoint left.} Let \(\alpha\) be a type. Let \(s\) and \(t\) be finite subsets of \(\alpha\). Then \(\operatorname{Disjoint}(s,\allowbreak{} t) \iff (\forall a\, :\, \alpha,\; a \in s \implies a \notin t)\).
\item[\paperlineanchor{premise:46696e7365742e656d7074795f737562736574}\paperlink{premise:46696e7365742e656d7074795f737562736574}{P19.}] \refstepcounter{paperpremise}\label{premise:46696e7365742e656d7074795f737562736574}\textbf{Empty subset.} Let \(\alpha\) be a type. Let \(s\) be a finite subset of \(\alpha\). Then \(\varnothing \subseteq s\).
\item[\paperlineanchor{premise:46696e7365742e65715f6f665f7375627365745f6f665f636172645f6c65}\paperlink{premise:46696e7365742e65715f6f665f7375627365745f6f665f636172645f6c65}{P20.}] \refstepcounter{paperpremise}\label{premise:46696e7365742e65715f6f665f7375627365745f6f665f636172645f6c65}\textbf{Equals of subset of cardinality at most.} Let \(\alpha\) be a type. Let \(s\) and \(t\) be finite subsets of \(\alpha\). Assume \(s \subseteq t\). Assume \(\left|t\right| \le \left|s\right|\). Then \(s = t\).
\item[\paperlineanchor{premise:46696e7365742e696e73744973416e746973796d6d537562736574}\paperlink{premise:46696e7365742e696e73744973416e746973796d6d537562736574}{P21.}] \refstepcounter{paperpremise}\label{premise:46696e7365742e696e73744973416e746973796d6d537562736574}\textbf{Inst is antisymm subset.} Let \(\alpha\) be a type. Then \(\operatorname{Is\,Antisymm}(\mathcal{P}_{\mathrm{fin}}(\alpha),\allowbreak{} \mathit{x1} \mapsto \mathit{x2} \mapsto \mathit{x1} \subseteq \mathit{x2})\).
\item[\paperlineanchor{premise:46696e7365742e696e737449734e6f6e7374726963745374726963744f7264657253756273657453537562736574}\paperlink{premise:46696e7365742e696e737449734e6f6e7374726963745374726963744f7264657253756273657453537562736574}{P22.}] \refstepcounter{paperpremise}\label{premise:46696e7365742e696e737449734e6f6e7374726963745374726963744f7264657253756273657453537562736574}\textbf{Inst is nonstrict strict order subset ssubset.} Let \(\alpha\) be a type. Then \(\operatorname{Is\,Nonstrict\,Strict\,Order}(\mathcal{P}_{\mathrm{fin}}(\alpha),\allowbreak{} \mathit{x1} \mapsto \mathit{x2} \mapsto \mathit{x1} \subseteq \mathit{x2},\allowbreak{} \mathit{x1} \mapsto \mathit{x2} \mapsto \mathit{x1} \subsetneq \mathit{x2})\).
\item[\paperlineanchor{premise:46696e7365742e696e737449735472616e73537562736574}\paperlink{premise:46696e7365742e696e737449735472616e73537562736574}{P23.}] \refstepcounter{paperpremise}\label{premise:46696e7365742e696e737449735472616e73537562736574}\textbf{Inst is trans subset.} Let \(\alpha\) be a type. Then \(\operatorname{Is\,Trans}(\mathcal{P}_{\mathrm{fin}}(\alpha),\allowbreak{} \mathit{x1} \mapsto \mathit{x2} \mapsto \mathit{x1} \subseteq \mathit{x2})\).
\item[\paperlineanchor{premise:46696e7365742e6d656d5f696e736572745f73656c66}\paperlink{premise:46696e7365742e6d656d5f696e736572745f73656c66}{P24.}] \refstepcounter{paperpremise}\label{premise:46696e7365742e6d656d5f696e736572745f73656c66}\textbf{Membership insert self.} Let \(\alpha\) be a type. Assume equality on \(\alpha\) is decidable. Let \(a\) be an element of \(\alpha\). Let \(s\) be a finite subset of \(\alpha\). Then \(a \in \{a\} \cup s\).
\item[\paperlineanchor{premise:46696e7365742e6d656d5f6d6170}\paperlink{premise:46696e7365742e6d656d5f6d6170}{P25.}] \refstepcounter{paperpremise}\label{premise:46696e7365742e6d656d5f6d6170}\textbf{Membership map.} Let \(\alpha\) and \(\beta\) be types. Let \(f\) be an element of \(\operatorname{Embedding}\left(\alpha, \beta\right)\). Let \(s\) be a finite subset of \(\alpha\). Let \(b\) be an element of \(\beta\). Then \(b \in \operatorname{map}(f,\allowbreak{} s) \iff (\exists a\, :\, \alpha,\; a \in s \land \operatorname{coe}(f,\allowbreak{} a) = b)\).
\item[\paperlineanchor{premise:46696e7365742e6e655f696e736572745f6f665f6e6f745f6d656d}\paperlink{premise:46696e7365742e6e655f696e736572745f6f665f6e6f745f6d656d}{P26.}] \refstepcounter{paperpremise}\label{premise:46696e7365742e6e655f696e736572745f6f665f6e6f745f6d656d}\textbf{Differs from insert of not membership.} Let \(\alpha\) be a type. Assume equality on \(\alpha\) is decidable. Let \(s\) and \(t\) be finite subsets of \(\alpha\). Let \(a\) be an element of \(\alpha\). Assume \(a \notin s\). Then \(s \ne \{a\} \cup t\).
\item[\paperlineanchor{premise:46696e7365742e73696e676c65746f6e5f6e655f656d707479}\paperlink{premise:46696e7365742e73696e676c65746f6e5f6e655f656d707479}{P27.}] \refstepcounter{paperpremise}\label{premise:46696e7365742e73696e676c65746f6e5f6e655f656d707479}\textbf{Singleton differs from empty.} Let \(\alpha\) be a type. Let \(a\) be an element of \(\alpha\). Then \(\{a\} \ne \varnothing\).
\item[\paperlineanchor{premise:46696e7365742e7375627365745f696e73657274}\paperlink{premise:46696e7365742e7375627365745f696e73657274}{P28.}] \refstepcounter{paperpremise}\label{premise:46696e7365742e7375627365745f696e73657274}\textbf{Subset insert.} Let \(\alpha\) be a type. Assume equality on \(\alpha\) is decidable. Let \(a\) be an element of \(\alpha\). Let \(s\) be a finite subset of \(\alpha\). Then \(s \subseteq \{a\} \cup s\).
\item[\paperlineanchor{premise:46756e6374696f6e2e456d62656464696e672e636f65466e5f6d6b}\paperlink{premise:46756e6374696f6e2e456d62656464696e672e636f65466e5f6d6b}{P29.}] \refstepcounter{paperpremise}\label{premise:46756e6374696f6e2e456d62656464696e672e636f65466e5f6d6b}\textbf{Coercion fn mk.} Let \(\alpha\) be a sort. Let \(\beta\) be a sort. Let \(f\) be a function of type \(\alpha \to \beta\). Let \(i\) be a proof of \(f\text{ is injective}\). Then \(\operatorname{coe}\left(\operatorname{mk}\left(f\right)\right) = f\).
\item[\paperlineanchor{premise:4861735375627365742e5375627365742e7472616e73}\paperlink{premise:4861735375627365742e5375627365742e7472616e73}{P30.}] \refstepcounter{paperpremise}\label{premise:4861735375627365742e5375627365742e7472616e73}\textbf{Trans.} Let \(\alpha\) be a type. Equip \(\alpha\) with a subset relation. Let \(\mathit{inst}_{1}\) be a proof of \(\operatorname{Is\,Trans}(\alpha,\allowbreak{} \mathit{x1} \mapsto \mathit{x2} \mapsto \mathit{x1} \subseteq \mathit{x2})\). Let \(a\), \(b\), and \(c\) be elements of \(\alpha\). Assume \(a \subseteq b\). Assume \(b \subseteq c\). Then \(a \subseteq c\).
\item[\paperlineanchor{premise:4966662e6e6f74}\paperlink{premise:4966662e6e6f74}{P31.}] \refstepcounter{paperpremise}\label{premise:4966662e6e6f74}\textbf{Not.} Let \(a\) and \(b\) be propositions. Assume \(a \iff b\). Then \(\neg a \iff \neg b\).
\item[\paperlineanchor{premise:497355707065725365742e6e6f745f746f705f6d656d}\paperlink{premise:497355707065725365742e6e6f745f746f705f6d656d}{P32.}] \refstepcounter{paperpremise}\label{premise:497355707065725365742e6e6f745f746f705f6d656d}\textbf{Not top membership.} Let \(\alpha\) be a type. Equip \(\alpha\) with a preorder. Let \(s\) be a subset of \(\alpha\). Fix a structure \(\mathit{inst}_{1}\) of type \(\operatorname{Order\,Top}\left(\alpha\right)\). Assume \(\operatorname{Is\,Upper\,Set}\left(s\right)\). Then \(\operatorname{top} \notin s \iff s = \varnothing\).
\item[\paperlineanchor{premise:4c452e6c652e7472616e735f6571}\paperlink{premise:4c452e6c652e7472616e735f6571}{P33.}] \refstepcounter{paperpremise}\label{premise:4c452e6c652e7472616e735f6571}\textbf{Trans equals.} Let \(\alpha\) be a type. Let \(a\), \(b\), and \(c\) be elements of \(\alpha\). Fix a structure \(\mathit{inst}\) of type \(\operatorname{LE}\left(\alpha\right)\). Assume \(a \le b\). Assume \(b = c\). Then \(a \le c\).
\item[\paperlineanchor{premise:4c452e6c652e7472616e735f657127}\paperlink{premise:4c452e6c652e7472616e735f657127}{P34.}] \refstepcounter{paperpremise}\label{premise:4c452e6c652e7472616e735f657127}\textbf{Trans eq'.} Let \(\alpha\) be a type. Equip \(\alpha\) with a preorder. Let \(a\), \(b\), and \(c\) be elements of \(\alpha\). Assume \(b \le c\). Assume \(a = b\). Then \(a \le c\).
\item[\paperlineanchor{premise:4d6174686c69622e5461637469632e507573684e65672e6e6f745f666f72616c6c5f6571}\paperlink{premise:4d6174686c69622e5461637469632e507573684e65672e6e6f745f666f72616c6c5f6571}{P35.}] \refstepcounter{paperpremise}\label{premise:4d6174686c69622e5461637469632e507573684e65672e6e6f745f666f72616c6c5f6571}\textbf{Not forall equals.} Let \(\alpha\) be a sort. Let \(s\) be a function of type \(\alpha \to \mathsf{Prop}\). Then \((\neg (\forall x\, :\, \alpha,\; s(x))) = (\exists x\, :\, \alpha,\; \neg s(x))\).
\item[\paperlineanchor{premise:4d6174686c69622e5461637469632e507573684e65672e6e6f745f696d706c6965735f6571}\paperlink{premise:4d6174686c69622e5461637469632e507573684e65672e6e6f745f696d706c6965735f6571}{P36.}] \refstepcounter{paperpremise}\label{premise:4d6174686c69622e5461637469632e507573684e65672e6e6f745f696d706c6965735f6571}\textbf{Not implies equals.} Let \(p\) and \(q\) be propositions. Then \(\left(\neg \left(p \implies q\right)\right) = \left(p \land \neg q\right)\).
\item[\paperlineanchor{premise:4d6174686c69622e5461637469632e507573684e65672e6e6f745f6f725f6571}\paperlink{premise:4d6174686c69622e5461637469632e507573684e65672e6e6f745f6f725f6571}{P37.}] \refstepcounter{paperpremise}\label{premise:4d6174686c69622e5461637469632e507573684e65672e6e6f745f6f725f6571}\textbf{Not or equals.} Let \(p\) and \(q\) be propositions. Then \(\left(\neg \left(p \lor q\right)\right) = \left(\neg p \land \neg q\right)\).
\item[\paperlineanchor{premise:4e61742e6c655f6469765f6966665f6d756c5f6c65}\paperlink{premise:4e61742e6c655f6469765f6966665f6d756c5f6c65}{P38.}] \refstepcounter{paperpremise}\label{premise:4e61742e6c655f6469765f6966665f6d756c5f6c65}\textbf{At most div if and only if product at most.} Let \(k\), \(x\), and \(y\) be natural numbers. Assume \(0 < k\). Then \(x \le y / k \iff x \cdot k \le y\).
\item[\paperlineanchor{premise:4e61742e6c655f6f665f6d756c5f6c655f6d756c5f6c656674}\paperlink{premise:4e61742e6c655f6f665f6d756c5f6c655f6d756c5f6c656674}{P39.}] \refstepcounter{paperpremise}\label{premise:4e61742e6c655f6f665f6d756c5f6c655f6d756c5f6c656674}\textbf{At most of product at most product left.} Let \(a\), \(b\) and \(c\) be natural numbers. If \(c \cdot a \le c \cdot b\) and \(0 < c\), then \(a \le b\).
\item[\paperlineanchor{premise:4e61742e6d756c5f636f6d6d}\paperlink{premise:4e61742e6d756c5f636f6d6d}{P40.}] \refstepcounter{paperpremise}\label{premise:4e61742e6d756c5f636f6d6d}\textbf{Multiplication commutativity.} Let \(n\) and \(m\) be natural numbers. Then \(n \cdot m = m \cdot n\).
\item[\paperlineanchor{premise:4e61742e74776f5f6d756c}\paperlink{premise:4e61742e74776f5f6d756c}{P41.}] \refstepcounter{paperpremise}\label{premise:4e61742e74776f5f6d756c}\textbf{Doubling by addition.} Let \(n\) be a natural number. Then \(2 \cdot n = n + n\).
\item[\paperlineanchor{premise:4e61742e74776f5f706f73}\paperlink{premise:4e61742e74776f5f706f73}{P42.}] \refstepcounter{paperpremise}\label{premise:4e61742e74776f5f706f73}\textbf{Two positivity.} \(0 < 2\).
\item[\paperlineanchor{premise:4e652e73796d6d}\paperlink{premise:4e652e73796d6d}{P43.}] \refstepcounter{paperpremise}\label{premise:4e652e73796d6d}\textbf{Symm.} Let \(\alpha\) be a sort. Let \(a\) and \(b\) be elements of \(\alpha\). Assume \(a \ne b\). Then \(b \ne a\).
\item[\paperlineanchor{premise:5365742e4e6f6e656d7074792e6e655f656d707479}\paperlink{premise:5365742e4e6f6e656d7074792e6e655f656d707479}{P44.}] \refstepcounter{paperpremise}\label{premise:5365742e4e6f6e656d7074792e6e655f656d707479}\textbf{Differs from empty.} Let \(\alpha\) be a type. Let \(s\) be a subset of \(\alpha\). Assume \(s\ne\varnothing\). Then \(s \ne \varnothing\).
\item[\paperlineanchor{premise:5365742e50616972776973652e6571}\paperlink{premise:5365742e50616972776973652e6571}{P45.}] \refstepcounter{paperpremise}\label{premise:5365742e50616972776973652e6571}\textbf{Equals.} Let \(\alpha\) be a type. Let \(r\) be a relation from \(\alpha\) to \(\alpha\). Let \(s\) be a subset of \(\alpha\). Let \(a\) and \(b\) be elements of \(\alpha\). Assume \(\operatorname{Pairwise}\left(s, r\right)\). Assume \(a \in s\). Assume \(b \in s\). Assume \(\neg r\left(a, b\right)\). Then \(a = b\).
\item[\paperlineanchor{premise:5365742e5375627365742e72666c}\paperlink{premise:5365742e5375627365742e72666c}{P46.}] \refstepcounter{paperpremise}\label{premise:5365742e5375627365742e72666c}\textbf{Rfl.} Let \(\alpha\) be a type. Let \(s\) be a subset of \(\alpha\). Then \(s \subseteq s\).
\item[\paperlineanchor{premise:5365742e53756273696e676c65746f6e2e65715f73696e676c65746f6e5f6f665f6d656d}\paperlink{premise:5365742e53756273696e676c65746f6e2e65715f73696e676c65746f6e5f6f665f6d656d}{P47.}] \refstepcounter{paperpremise}\label{premise:5365742e53756273696e676c65746f6e2e65715f73696e676c65746f6e5f6f665f6d656d}\textbf{Equals singleton of membership.} Let \(\alpha\) be a type. Let \(s\) be a subset of \(\alpha\). Assume \(\operatorname{Set\mathord{:}Subsingleton}\left(s\right)\). Let \(x\) be an element of \(\alpha\). Assume \(x \in s\). Then \(s = \{x\}\).
\item[\paperlineanchor{premise:5365742e53756273696e676c65746f6e2e7061697277697365}\paperlink{premise:5365742e53756273696e676c65746f6e2e7061697277697365}{P48.}] \refstepcounter{paperpremise}\label{premise:5365742e53756273696e676c65746f6e2e7061697277697365}\textbf{Pairwise.} Let \(\alpha\) be a type. Let \(s\) be a subset of \(\alpha\). Assume \(\operatorname{Set\mathord{:}Subsingleton}\left(s\right)\). Let \(r\) be a relation from \(\alpha\) to \(\alpha\). Then \(\operatorname{Pairwise}\left(s, r\right)\).
\item[\paperlineanchor{premise:5365742e65715f73696e676c65746f6e5f6966665f756e697175655f6d656d}\paperlink{premise:5365742e65715f73696e676c65746f6e5f6966665f756e697175655f6d656d}{P49.}] \refstepcounter{paperpremise}\label{premise:5365742e65715f73696e676c65746f6e5f6966665f756e697175655f6d656d}\textbf{Equals singleton if and only if unique membership.} Let \(\alpha\) be a type. Let \(s\) be a subset of \(\alpha\). Let \(a\) be an element of \(\alpha\). Then \(s = \{a\} \iff a \in s \land \left(\forall x\, :\, \alpha,\; x \in s \implies x = a\right)\).
\item[\paperlineanchor{premise:5365742e65715f756e69765f6f665f666f72616c6c}\paperlink{premise:5365742e65715f756e69765f6f665f666f72616c6c}{P50.}] \refstepcounter{paperpremise}\label{premise:5365742e65715f756e69765f6f665f666f72616c6c}\textbf{Equals whole space of forall.} Let \(\alpha\) be a type. Let \(s\) be a subset of \(\alpha\). Assume \(\forall x\, :\, \alpha,\; x \in s\). Then \(s = \alpha\).
\item[\paperlineanchor{premise:5365742e696d6167655f65715f707265696d6167655f6f665f696e7665727365}\paperlink{premise:5365742e696d6167655f65715f707265696d6167655f6f665f696e7665727365}{P51.}] \refstepcounter{paperpremise}\label{premise:5365742e696d6167655f65715f707265696d6167655f6f665f696e7665727365}\textbf{Image equals preimage of inverse.} Let \(\alpha\) and \(\beta\) be types. Let \(f\) be a function of type \(\alpha \to \beta\). Let \(g\) be a function of type \(\beta \to \alpha\). Assume \(\operatorname{Left\,Inverse}\left(g, f\right)\). Assume \(\operatorname{Right\,Inverse}\left(g, f\right)\). Then \(\operatorname{image}\left(f\right) = \operatorname{preimage}\left(g\right)\).
\item[\paperlineanchor{premise:5365742e6d656d5f696e73657274}\paperlink{premise:5365742e6d656d5f696e73657274}{P52.}] \refstepcounter{paperpremise}\label{premise:5365742e6d656d5f696e73657274}\textbf{Membership insert.} Let \(\alpha\) be a type. Let \(x\) be an element of \(\alpha\). Let \(s\) be a subset of \(\alpha\). Then \(x \in \{x\} \cup s\).
\item[\paperlineanchor{premise:5365742e6d656d5f696e736572745f6f665f6d656d}\paperlink{premise:5365742e6d656d5f696e736572745f6f665f6d656d}{P53.}] \refstepcounter{paperpremise}\label{premise:5365742e6d656d5f696e736572745f6f665f6d656d}\textbf{Membership insert of membership.} Let \(\alpha\) be a type. Let \(x\) be an element of \(\alpha\). Let \(s\) be a subset of \(\alpha\). Let \(y\) be an element of \(\alpha\). Assume \(x \in s\). Then \(x \in \{y\} \cup s\).
\item[\paperlineanchor{premise:5365742e6d656d5f707265696d616765}\paperlink{premise:5365742e6d656d5f707265696d616765}{P54.}] \refstepcounter{paperpremise}\label{premise:5365742e6d656d5f707265696d616765}\textbf{Membership preimage.} Let \(\alpha\) and \(\beta\) be types. Let \(f\) be a function of type \(\alpha \to \beta\). Let \(s\) be a subset of \(\beta\). Let \(a\) be an element of \(\alpha\). Then \(a \in f^{-1}\left[s\right] \iff f\left(a\right) \in s\).
\item[\paperlineanchor{premise:5365742e6d656d5f73696e676c65746f6e5f696666}\paperlink{premise:5365742e6d656d5f73696e676c65746f6e5f696666}{P55.}] \refstepcounter{paperpremise}\label{premise:5365742e6d656d5f73696e676c65746f6e5f696666}\textbf{Membership singleton if and only if.} Let \(\alpha\) be a type. Let \(a\) and \(b\) be elements of \(\alpha\). Then \(a \in \{b\} \iff a = b\).
\item[\paperlineanchor{premise:5365742e6d656d5f756e696f6e}\paperlink{premise:5365742e6d656d5f756e696f6e}{P56.}] \refstepcounter{paperpremise}\label{premise:5365742e6d656d5f756e696f6e}\textbf{Membership union.} Let \(\alpha\) be a type. Let \(x\) be an element of \(\alpha\). Let \(a\) and \(b\) be subsets of \(\alpha\). Then \(x \in a \cup b \iff x \in a \lor x \in b\).
\item[\paperlineanchor{premise:5365742e6e6f6e656d7074795f6966665f6e655f656d707479}\paperlink{premise:5365742e6e6f6e656d7074795f6966665f6e655f656d707479}{P57.}] \refstepcounter{paperpremise}\label{premise:5365742e6e6f6e656d7074795f6966665f6e655f656d707479}\textbf{Nonempty if and only if differs from empty.} Let \(\alpha\) be a type. Let \(s\) be a subset of \(\alpha\). Then \(s\ne\varnothing \iff s \ne \varnothing\).
\item[\paperlineanchor{premise:5365742e6e6f745f6469736a6f696e745f696666}\paperlink{premise:5365742e6e6f745f6469736a6f696e745f696666}{P58.}] \refstepcounter{paperpremise}\label{premise:5365742e6e6f745f6469736a6f696e745f696666}\textbf{Not disjoint if and only if.} Let \(\alpha\) be a type. Let \(s\) and \(t\) be subsets of \(\alpha\). Then \(\neg \operatorname{Disjoint}(s,\allowbreak{} t) \iff (\exists x\, :\, \alpha,\; x \in s \land x \in t)\).
\item[\paperlineanchor{premise:5365742e73696e676c65746f6e5f6e6f6e656d707479}\paperlink{premise:5365742e73696e676c65746f6e5f6e6f6e656d707479}{P59.}] \refstepcounter{paperpremise}\label{premise:5365742e73696e676c65746f6e5f6e6f6e656d707479}\textbf{Singleton nonempty.} Let \(\alpha\) be a type. Let \(a\) be an element of \(\alpha\). Then \(\{a\}\ne\varnothing\).
\item[\paperlineanchor{premise:5365742e7375627365745f696e73657274}\paperlink{premise:5365742e7375627365745f696e73657274}{P60.}] \refstepcounter{paperpremise}\label{premise:5365742e7375627365745f696e73657274}\textbf{Subset insert.} Let \(\alpha\) be a type. Let \(x\) be an element of \(\alpha\). Let \(s\) be a subset of \(\alpha\). Then \(s \subseteq \{x\} \cup s\).
\item[\paperlineanchor{premise:5365742e73756273696e676c65746f6e5f6f665f73756273696e676c65746f6e}\paperlink{premise:5365742e73756273696e676c65746f6e5f6f665f73756273696e676c65746f6e}{P61.}] \refstepcounter{paperpremise}\label{premise:5365742e73756273696e676c65746f6e5f6f665f73756273696e676c65746f6e}\textbf{Subsingleton of subsingleton.} Let \(\alpha\) be a type. Assume \(\alpha\text{ has at most one element}\). Let \(s\) be a subset of \(\alpha\). Then \(\operatorname{Set\mathord{:}Subsingleton}\left(s\right)\).
\item[\paperlineanchor{premise:53756273696e676c65746f6e2e656c696d}\paperlink{premise:53756273696e676c65746f6e2e656c696d}{P62.}] \refstepcounter{paperpremise}\label{premise:53756273696e676c65746f6e2e656c696d}\textbf{Elim.} Let \(\alpha\) be a sort. Assume \(\alpha\text{ has at most one element}\). Let \(a\) and \(b\) be elements of \(\alpha\). Then \(a = b\).
\item[\paperlineanchor{premise:616e645f6966665f7269676874}\paperlink{premise:616e645f6966665f7269676874}{P63.}] \refstepcounter{paperpremise}\label{premise:616e645f6966665f7269676874}\textbf{And if and only if right.} Let \(a\) and \(b\) be propositions. Assume \(a\). Then \(a \land b \iff b\).
\item[\paperlineanchor{premise:636f6d706c5f626f74}\paperlink{premise:636f6d706c5f626f74}{P64.}] \refstepcounter{paperpremise}\label{premise:636f6d706c5f626f74}\textbf{Compl bot.} Let \(\alpha\) be a type. Fix a structure \(\mathit{inst}\) of type \(\operatorname{Heyting\,Algebra}\left(\alpha\right)\). Then \(\operatorname{compl}\left(\operatorname{bot}\right) = \operatorname{top}\).
\item[\paperlineanchor{premise:636f6d706c5f636f6d706c}\paperlink{premise:636f6d706c5f636f6d706c}{P65.}] \refstepcounter{paperpremise}\label{premise:636f6d706c5f636f6d706c}\textbf{Compl compl.} Let \(\alpha\) be a type. Fix a structure \(\mathit{inst}\) of type \(\operatorname{Boolean\,Algebra}\left(\alpha\right)\). Let \(x\) be an element of \(\alpha\). Then \(\operatorname{compl}\left(\operatorname{compl}\left(x\right)\right) = x\).
\item[\paperlineanchor{premise:636f6d706c5f696e6a656374697665}\paperlink{premise:636f6d706c5f696e6a656374697665}{P66.}] \refstepcounter{paperpremise}\label{premise:636f6d706c5f696e6a656374697665}\textbf{Compl injective.} Let \(\alpha\) be a type. Fix a structure \(\mathit{inst}\) of type \(\operatorname{Boolean\,Algebra}\left(\alpha\right)\). Then \(\operatorname{compl}\text{ is injective}\).
\item[\paperlineanchor{premise:6469736a6f696e745f626f745f6c656674}\paperlink{premise:6469736a6f696e745f626f745f6c656674}{P67.}] \refstepcounter{paperpremise}\label{premise:6469736a6f696e745f626f745f6c656674}\textbf{Disjoint bot left.} Let \(\alpha\) be a type. Equip \(\alpha\) with a partial order. Fix a structure \(\mathit{inst}_{1}\) of type \(\operatorname{Order\,Bot}\left(\alpha\right)\). Let \(a\) be an element of \(\alpha\). Then \(\operatorname{Disjoint}\left(\operatorname{bot}, a\right)\).
\item[\paperlineanchor{premise:6469736a6f696e745f636f6d706c5f6c656674}\paperlink{premise:6469736a6f696e745f636f6d706c5f6c656674}{P68.}] \refstepcounter{paperpremise}\label{premise:6469736a6f696e745f636f6d706c5f6c656674}\textbf{Disjoint compl left.} Let \(\alpha\) be a type. Fix a structure \(\mathit{inst}\) of type \(\operatorname{Heyting\,Algebra}\left(\alpha\right)\). Let \(a\) be an element of \(\alpha\). Then \(\operatorname{Disjoint}\left(\operatorname{compl}\left(a\right), a\right)\).
\item[\paperlineanchor{premise:6469736a6f696e745f636f6d706c5f7269676874}\paperlink{premise:6469736a6f696e745f636f6d706c5f7269676874}{P69.}] \refstepcounter{paperpremise}\label{premise:6469736a6f696e745f636f6d706c5f7269676874}\textbf{Disjoint compl right.} Let \(\alpha\) be a type. Fix a structure \(\mathit{inst}\) of type \(\operatorname{Heyting\,Algebra}\left(\alpha\right)\). Let \(a\) be an element of \(\alpha\). Then \(\operatorname{Disjoint}\left(a, \operatorname{compl}\left(a\right)\right)\).
\item[\paperlineanchor{premise:6469736a6f696e745f73656c66}\paperlink{premise:6469736a6f696e745f73656c66}{P70.}] \refstepcounter{paperpremise}\label{premise:6469736a6f696e745f73656c66}\textbf{Disjoint self.} Let \(\alpha\) be a type. Equip \(\alpha\) with a partial order. Fix a structure \(\mathit{inst}_{1}\) of type \(\operatorname{Order\,Bot}\left(\alpha\right)\). Let \(a\) be an element of \(\alpha\). Then \(\operatorname{Disjoint}\left(a, a\right) \iff a = \operatorname{bot}\).
\item[\paperlineanchor{premise:65715f626f745f6d6f6e6f}\paperlink{premise:65715f626f745f6d6f6e6f}{P71.}] \refstepcounter{paperpremise}\label{premise:65715f626f745f6d6f6e6f}\textbf{Equals bot mono.} Let \(\alpha\) be a type. Equip \(\alpha\) with a partial order. Fix a structure \(\mathit{inst}_{1}\) of type \(\operatorname{Order\,Bot}\left(\alpha\right)\). Let \(a\) and \(b\) be elements of \(\alpha\). Assume \(a \le b\). Assume \(b = \operatorname{bot}\). Then \(a = \operatorname{bot}\).
\item[\paperlineanchor{premise:666f72616c6c5f636f6e6772}\paperlink{premise:666f72616c6c5f636f6e6772}{P72.}] \refstepcounter{paperpremise}\label{premise:666f72616c6c5f636f6e6772}\textbf{Forall congruence.} Let \(\alpha\) be a sort. Let \(p\) and \(q\) be functions of type \(\alpha \to \mathsf{Prop}\). Assume \(\forall a\, :\, \alpha,\; p\left(a\right) = q\left(a\right)\). Then \((\forall a\, :\, \alpha,\; p(a)) = (\forall a\, :\, \alpha,\; q(a))\).
\item[\paperlineanchor{premise:666f72616c6c5f6571}\paperlink{premise:666f72616c6c5f6571}{P73.}] \refstepcounter{paperpremise}\label{premise:666f72616c6c5f6571}\textbf{Forall equals.} Let \(\alpha\) be a sort. Let \(p\) be a function of type \(\alpha \to \mathsf{Prop}\). Let \(a^{'}\) be an element of \(\alpha\). Then \(\left(\forall a\, :\, \alpha,\; a = a^{'} \implies p\left(a\right)\right) \iff p\left(a^{'}\right)\).
\item[\paperlineanchor{premise:66756e657874}\paperlink{premise:66756e657874}{P74.}] \refstepcounter{paperpremise}\label{premise:66756e657874}\textbf{Function extensionality.} Let \(\alpha\) be a sort. Let \(\beta\) be a function of type \(\alpha \to \mathsf{Sort}\). Let \(f\) be a function of type \(\Pi_{x\, :\, \alpha}\, \beta\left(x\right)\). Let \(g\) be a function of type \(\Pi_{x\, :\, \alpha}\, \beta\left(x\right)\). Assume \(\forall x\, :\, \alpha,\; f\left(x\right) = g\left(x\right)\). Then \(f = g\).
\item[\paperlineanchor{premise:6966665f73656c66}\paperlink{premise:6966665f73656c66}{P75.}] \refstepcounter{paperpremise}\label{premise:6966665f73656c66}\textbf{If and only if self.} Let \(p\) be a proposition. Then \(\left(p \iff p\right) = \mathrm{True}\).
\item[\paperlineanchor{premise:696d706c6965735f636f6e6772}\paperlink{premise:696d706c6965735f636f6e6772}{P76.}] \refstepcounter{paperpremise}\label{premise:696d706c6965735f636f6e6772}\textbf{Implies congruence.} Let \(p_{1}\) and \(p_{2}\) be sorts. Let \(q_{1}\) and \(q_{2}\) be sorts. Assume \(p_{1} = p_{2}\). Assume \(q_{1} = q_{2}\). Then \(\left(p_{1} \to q_{1}\right) = \left(p_{2} \to q_{2}\right)\).
\item[\paperlineanchor{premise:6d74}\paperlink{premise:6d74}{P77.}] \refstepcounter{paperpremise}\label{premise:6d74}\textbf{Mt.} Let \(a\) and \(b\) be propositions. Assume \(a \implies b\). Assume \(\neg b\). Then \(\neg a\).
\item[\paperlineanchor{premise:6e655f6f665f6d656d5f6f665f6e6f745f6d656d}\paperlink{premise:6e655f6f665f6d656d5f6f665f6e6f745f6d656d}{P78.}] \refstepcounter{paperpremise}\label{premise:6e655f6f665f6d656d5f6f665f6e6f745f6d656d}\textbf{Differs from of membership of not membership.} Let \(\alpha\) and \(\beta\) be types. Fix a structure \(\mathit{inst}\) of type \(\operatorname{Membership}\left(\alpha, \beta\right)\). Let \(s\) be an element of \(\beta\). Let \(a\) and \(b\) be elements of \(\alpha\). Assume \(a \in s\). Assume \(b \notin s\). Then \(a \ne b\).
\item[\paperlineanchor{premise:6e6f745f696d705f636f6d6d}\paperlink{premise:6e6f745f696d705f636f6d6d}{P79.}] \refstepcounter{paperpremise}\label{premise:6e6f745f696d705f636f6d6d}\textbf{Not imp commutation.} Let \(a\) and \(b\) be propositions. Then \(\neg a \implies b \iff \neg b \implies a\).
\item[\paperlineanchor{premise:6e6f745f6e655f696666}\paperlink{premise:6e6f745f6e655f696666}{P80.}] \refstepcounter{paperpremise}\label{premise:6e6f745f6e655f696666}\textbf{Not differs from if and only if.} Let \(\alpha\) be a sort. Let \(a\) and \(b\) be elements of \(\alpha\). Then \(\neg \left(a \ne b\right) \iff a = b\).
\item[\paperlineanchor{premise:737375627365745f6966665f7375627365745f6e65}\paperlink{premise:737375627365745f6966665f7375627365745f6e65}{P81.}] \refstepcounter{paperpremise}\label{premise:737375627365745f6966665f7375627365745f6e65}\textbf{Ssubset if and only if subset differs from.} Let \(\alpha\) be a type. Equip \(\alpha\) with a subset relation. Fix a structure \(\mathit{inst}_{1}\) of type \(\operatorname{Has\,SSubset}\left(\alpha\right)\). Let \(\mathit{inst}_{2}\) be a proof of \(\operatorname{Is\,Nonstrict\,Strict\,Order}(\alpha,\allowbreak{} \mathit{x1} \mapsto \mathit{x2} \mapsto \mathit{x1} \subseteq \mathit{x2},\allowbreak{} \mathit{x1} \mapsto \mathit{x2} \mapsto \mathit{x1} \subsetneq \mathit{x2})\). Let \(a\) and \(b\) be elements of \(\alpha\). Let \(\mathit{inst}_{3}\) be a proof of \(\operatorname{Is\,Antisymm}(\alpha,\allowbreak{} \mathit{x1} \mapsto \mathit{x2} \mapsto \mathit{x1} \subseteq \mathit{x2})\). Then \(a \subsetneq b \iff a \subseteq b \land a \ne b\).
\item[\paperlineanchor{premise:746f705f6e655f626f74}\paperlink{premise:746f705f6e655f626f74}{P82.}] \refstepcounter{paperpremise}\label{premise:746f705f6e655f626f74}\textbf{Top differs from bot.} Let \(\alpha\) be a type. Equip \(\alpha\) with a partial order. Fix a structure \(\mathit{inst}_{1}\) of type \(\operatorname{Bounded\,Order}\left(\alpha\right)\). Let \(\mathit{inst}_{2}\) be a proof of \(\operatorname{Nontrivial}\left(\alpha\right)\). Then \(\operatorname{top} \ne \operatorname{bot}\).
\end{description}

\section{Semilattice inf}
\begin{lemma}[{Mono}]\label{theorem:1}
Let \(\alpha\) be a type. Fix a structure \(\mathit{inst}\) of type \(\operatorname{Semilattice\,Inf}\left(\alpha\right)\). Fix a structure \(\mathit{inst}_{1}\) of type \(\operatorname{Order\,Bot}\left(\alpha\right)\). Let \(s\) and \(t\) be subsets of \(\alpha\). Assume \(t \subseteq s\). Assume \(\operatorname{Intersecting}\left(s\right)\). Then \(\operatorname{Intersecting}\left(t\right)\).
\end{lemma}

\begin{proof}
Fix \(x\) in \(\alpha\).
Assume \(x \in t\).
Fix \(x_{1}\) in \(\alpha\).
Assume \(x_{1} \in t\).
Use the hypothesis \(x \in t\).
Applying the hypothesis \(t \subseteq s\), we obtain \(x \in s\).
Use the hypothesis \(x_{1} \in t\).
Applying the hypothesis \(t \subseteq s\), we obtain \(x_{1} \in s\).
Applying the hypothesis \(\operatorname{Intersecting}\left(s\right)\), we obtain \(\neg \operatorname{Disjoint}\left(x, x_{1}\right)\).
\end{proof}

\begin{lemma}[{Not bot membership}]\label{theorem:2}
Let \(\alpha\) be a type. Fix a structure \(\mathit{inst}\) of type \(\operatorname{Semilattice\,Inf}\left(\alpha\right)\). Fix a structure \(\mathit{inst}_{1}\) of type \(\operatorname{Order\,Bot}\left(\alpha\right)\). Let \(s\) be a subset of \(\alpha\). Assume \(\operatorname{Intersecting}\left(s\right)\). Then \(\operatorname{bot} \notin s\).
\end{lemma}

\begin{proof}
Assume \(\operatorname{bot} \in s\).
Use the hypothesis \(\operatorname{bot} \in s\).
Use the hypothesis \(\operatorname{bot} \in s\).
By \paperlink{premise:6469736a6f696e745f626f745f6c656674}{Prerequisite~P\ref{premise:6469736a6f696e745f626f745f6c656674}}, we obtain \(\operatorname{Disjoint}\left(\operatorname{bot}, \operatorname{bot}\right)\).
Applying the hypothesis \(\operatorname{Intersecting}\left(s\right)\), we obtain \(\mathrm{False}\).
\end{proof}

\begin{lemma}[{Differs from bot}]\label{theorem:3}
Let \(\alpha\) be a type. Fix a structure \(\mathit{inst}\) of type \(\operatorname{Semilattice\,Inf}\left(\alpha\right)\). Fix a structure \(\mathit{inst}_{1}\) of type \(\operatorname{Order\,Bot}\left(\alpha\right)\). Let \(s\) be a subset of \(\alpha\). Let \(a\) be an element of \(\alpha\). Assume \(\operatorname{Intersecting}\left(s\right)\). Assume \(a \in s\). Then \(a \ne \operatorname{bot}\).
\end{lemma}

\begin{proof}
Use the hypothesis \(a \in s\).
By \paperlink{theorem:2}{Lemma~\ref{theorem:2}}, we obtain \(\operatorname{bot} \notin s\).
By \paperlink{premise:6e655f6f665f6d656d5f6f665f6e6f745f6d656d}{Prerequisite~P\ref{premise:6e655f6f665f6d656d5f6f665f6e6f745f6d656d}} with \(\beta = \mathcal{P}(\alpha)\), we obtain \(a \ne \operatorname{bot}\).
\end{proof}

\begin{theorem}[{Intersecting empty}]\label{theorem:4}
Let \(\alpha\) be a type. Fix a structure \(\mathit{inst}\) of type \(\operatorname{Semilattice\,Inf}\left(\alpha\right)\). Fix a structure \(\mathit{inst}_{1}\) of type \(\operatorname{Order\,Bot}\left(\alpha\right)\). Then \(\operatorname{Intersecting}\left(\varnothing\right)\).
\end{theorem}

\begin{proof}
Fix \(x\) in \(\alpha\).
The contradiction rules out this case. Hence \(\mathrm{False} \implies (\forall b\, :\, \alpha,\; b \in \varnothing \implies \neg \operatorname{Disjoint}(x,\allowbreak{} b))\).
\end{proof}

\begin{lemma}[{Intersecting singleton}]\label{theorem:5}
Let \(\alpha\) be a type. Fix a structure \(\mathit{inst}\) of type \(\operatorname{Semilattice\,Inf}\left(\alpha\right)\). Fix a structure \(\mathit{inst}_{1}\) of type \(\operatorname{Order\,Bot}\left(\alpha\right)\). Let \(a\) be an element of \(\alpha\). Then \(\operatorname{Intersecting}\left(\{a\}\right) \iff a \ne \operatorname{bot}\).
\end{lemma}

\begin{proof}
The substitution uses the following argument.
\begin{proofstep}
Fix \(a_{1}\) in \(\alpha\).
By \paperlink{premise:5365742e6d656d5f73696e676c65746f6e5f696666}{Prerequisite~P\ref{premise:5365742e6d656d5f73696e676c65746f6e5f696666}}, we obtain \(a_{1} \in \{a\} \iff a_{1} = a\).
By propositional extensionality, we obtain \(\left(a_{1} \in \{a\}\right) = \left(a_{1} = a\right)\).
For the required implication, argue as follows.
\begin{proofstep}
Fix \(b\) in \(\alpha\).
By \paperlink{premise:5365742e6d656d5f73696e676c65746f6e5f696666}{Prerequisite~P\ref{premise:5365742e6d656d5f73696e676c65746f6e5f696666}}, we obtain \(b \in \{a\} \iff b = a\).
By propositional extensionality, we obtain \(\left(b \in \{a\}\right) = \left(b = a\right)\).
By definition, we obtain \((\neg \operatorname{Disjoint}(a_{1},\allowbreak{} b)) = (\neg \operatorname{Disjoint}(a_{1},\allowbreak{} b))\).
By \paperlink{premise:696d706c6965735f636f6e6772}{Prerequisite~P\ref{premise:696d706c6965735f636f6e6772}} with \(p_{1} = b \in \{a\}\) and \(p_{2} = b = a\) and \(q_{1} = \neg \operatorname{Disjoint}\left(a_{1}, b\right)\) and \(q_{2} = \neg \operatorname{Disjoint}\left(a_{1}, b\right)\), we obtain \((b \in \{a\} \implies \neg \operatorname{Disjoint}(a_{1},\allowbreak{} b)) = (b = a \implies \neg \operatorname{Disjoint}(a_{1},\allowbreak{} b))\).
\end{proofstep}
By \paperlink{premise:666f72616c6c5f636f6e6772}{Prerequisite~P\ref{premise:666f72616c6c5f636f6e6772}}, we obtain \((\forall a_{2}\, :\, \alpha,\; a_{2} \in \{a\} \implies \neg \operatorname{Disjoint}(a_{1},\allowbreak{} a_{2})) = (\forall a_{2}\, :\, \alpha,\; a_{2} = a \implies \neg \operatorname{Disjoint}(a_{1},\allowbreak{} a_{2}))\).
By \paperlink{premise:666f72616c6c5f6571}{Prerequisite~P\ref{premise:666f72616c6c5f6571}}, we obtain \((\forall a_{2}\, :\, \alpha,\; a_{2} = a \implies (a_{4} \mapsto \neg \operatorname{Disjoint}(a_{1},\allowbreak{} a_{4}))(a_{2})) \iff (a_{2} \mapsto \neg \operatorname{Disjoint}(a_{1},\allowbreak{} a_{2}))(a)\).
By propositional extensionality, we obtain \((\forall a_{2}\, :\, \alpha,\; a_{2} = a \implies \neg \operatorname{Disjoint}(a_{1},\allowbreak{} a_{2})) = (\neg \operatorname{Disjoint}(a_{1},\allowbreak{} a))\).
Chaining the equalities, we obtain \((\forall a_{2}\, :\, \alpha,\; a_{2} \in \{a\} \implies \neg \operatorname{Disjoint}(a_{1},\allowbreak{} a_{2})) = (\neg \operatorname{Disjoint}(a_{1},\allowbreak{} a))\).
By \paperlink{premise:696d706c6965735f636f6e6772}{Prerequisite~P\ref{premise:696d706c6965735f636f6e6772}} with \(p_{1} = a_{1} \in \{a\}\) and \(p_{2} = a_{1} = a\) and \(q_{1} = \forall a_{2}\, :\, \alpha,\; a_{2} \in \{a\} \implies \neg \operatorname{Disjoint}(a_{1},\allowbreak{} a_{2})\) and \(q_{2} = \neg \operatorname{Disjoint}\left(a_{1}, a\right)\), we obtain \((a_{1} \in \{a\} \implies (\forall a_{3}\, :\, \alpha,\; a_{3} \in \{a\} \implies \neg \operatorname{Disjoint}(a_{1},\allowbreak{} a_{3}))) = (a_{1} = a \implies \neg \operatorname{Disjoint}(a_{1},\allowbreak{} a))\).
\end{proofstep}
Rewriting with \paperlink{premise:666f72616c6c5f636f6e6772}{Prerequisite~P\ref{premise:666f72616c6c5f636f6e6772}}, \paperlink{premise:666f72616c6c5f6571}{Prerequisite~P\ref{premise:666f72616c6c5f6571}}, \paperlink{premise:6469736a6f696e745f73656c66}{Prerequisite~P\ref{premise:6469736a6f696e745f73656c66}}, \paperlink{premise:6966665f73656c66}{Prerequisite~P\ref{premise:6966665f73656c66}} and the identity established above gives \(\operatorname{Intersecting}\left(\{a\}\right) \iff a \ne \operatorname{bot}\).
\end{proof}

\begin{lemma}[{Insert}]\label{theorem:6}
Let \(\alpha\) be a type. Fix a structure \(\mathit{inst}\) of type \(\operatorname{Semilattice\,Inf}\left(\alpha\right)\). Fix a structure \(\mathit{inst}_{1}\) of type \(\operatorname{Order\,Bot}\left(\alpha\right)\). Let \(s\) be a subset of \(\alpha\). Let \(a\) be an element of \(\alpha\). Assume \(\operatorname{Intersecting}\left(s\right)\). Assume \(a \ne \operatorname{bot}\). Assume \(\forall b\, :\, \alpha,\; b \in s \implies \neg \operatorname{Disjoint}\left(a, b\right)\). Then \(\operatorname{Intersecting}\left(\{a\} \cup s\right)\).
\end{lemma}

\begin{proof}
Fix \(b\) in \(\alpha\).
Assume \(b \in \{a\} \cup s\).
Use the hypothesis \(b \in \{a\} \cup s\).
Consider the following case.
\begin{proofstep}
Assume \(b = a\).
For the required implication, argue as follows.
\begin{proofstep}
Assume \(b \ne \operatorname{bot}\).
Assume \(\forall b_{1}\, :\, \alpha,\; b_{1} \in s \implies \neg \operatorname{Disjoint}\left(b, b_{1}\right)\).
Fix \(c\) in \(\alpha\).
Assume \(c \in \{b\} \cup s\).
Use the hypothesis \(c \in \{b\} \cup s\).
Consider the following case.
\begin{proofstep}
Assume \(c = b\).
For the required implication, argue as follows.
\begin{proofstep}
Assume \(c \ne \operatorname{bot}\).
Assume \(\forall b_{1}\, :\, \alpha,\; b_{1} \in s \implies \neg \operatorname{Disjoint}\left(c, b_{1}\right)\).
Use the hypothesis \(c \ne \operatorname{bot}\), which simplifies to \(\neg \left(c = \operatorname{bot}\right)\).
Rewriting with \paperlink{premise:6469736a6f696e745f73656c66}{Prerequisite~P\ref{premise:6469736a6f696e745f73656c66}} in the preceding result gives \(\neg \operatorname{Disjoint}\left(c, c\right)\).
\end{proofstep}
Use the hypothesis \(c = b\).
Use the hypothesis \(b \ne \operatorname{bot}\).
Use the hypothesis \(\forall b_{1}\, :\, \alpha,\; b_{1} \in s \implies \neg \operatorname{Disjoint}\left(b, b_{1}\right)\).
Substituting equals, we obtain \(\neg \operatorname{Disjoint}\left(b, c\right)\).
\end{proofstep}
Consider the following case.
\begin{proofstep}
Assume \(c \in s\).
Use the hypothesis \(c \in s\).
Applying the hypothesis \(\forall b_{1}\, :\, \alpha,\; b_{1} \in s \implies \neg \operatorname{Disjoint}\left(b, b_{1}\right)\), we obtain \(\neg \operatorname{Disjoint}\left(b, c\right)\).
\end{proofstep}
The two alternatives exhaust the cases. Hence \(\neg \operatorname{Disjoint}\left(b, c\right)\).
\end{proofstep}
Use the hypothesis \(b = a\).
Use the hypothesis \(a \ne \operatorname{bot}\).
Use the hypothesis \(\forall b\, :\, \alpha,\; b \in s \implies \neg \operatorname{Disjoint}\left(a, b\right)\).
Substituting equals, we obtain \(\forall b_{1}\, :\, \alpha,\; b_{1} \in \{a\} \cup s \implies \neg \operatorname{Disjoint}(b,\allowbreak{} b_{1})\).
\end{proofstep}
Consider the following case.
\begin{proofstep}
Assume \(b \in s\).
Fix \(c\) in \(\alpha\).
Assume \(c \in \{a\} \cup s\).
Use the hypothesis \(c \in \{a\} \cup s\).
Consider the following case.
\begin{proofstep}
Assume \(c = a\).
For the required implication, argue as follows.
\begin{proofstep}
Assume \(c \ne \operatorname{bot}\).
Assume \(\forall b_{1}\, :\, \alpha,\; b_{1} \in s \implies \neg \operatorname{Disjoint}\left(c, b_{1}\right)\).
Assume \(\operatorname{Disjoint}\left(b, c\right)\).
Use the hypothesis \(b \in s\).
Use the hypothesis \(\operatorname{Disjoint}\left(b, c\right)\).
By \paperlink{premise:4469736a6f696e742e73796d6d}{Prerequisite~P\ref{premise:4469736a6f696e742e73796d6d}}, we obtain \(\operatorname{Disjoint}\left(c, b\right)\).
Applying the hypothesis \(\forall b_{1}\, :\, \alpha,\; b_{1} \in s \implies \neg \operatorname{Disjoint}\left(c, b_{1}\right)\), we obtain \(\mathrm{False}\).
\end{proofstep}
Use the hypothesis \(c = a\).
Use the hypothesis \(a \ne \operatorname{bot}\).
Use the hypothesis \(\forall b\, :\, \alpha,\; b \in s \implies \neg \operatorname{Disjoint}\left(a, b\right)\).
Substituting equals, we obtain \(\neg \operatorname{Disjoint}\left(b, c\right)\).
\end{proofstep}
Consider the following case.
\begin{proofstep}
Assume \(c \in s\).
Use the hypothesis \(b \in s\).
Use the hypothesis \(c \in s\).
Applying the hypothesis \(\operatorname{Intersecting}\left(s\right)\), we obtain \(\neg \operatorname{Disjoint}\left(b, c\right)\).
\end{proofstep}
The two alternatives exhaust the cases. Hence \(\neg \operatorname{Disjoint}\left(b, c\right)\).
\end{proofstep}
The two alternatives exhaust the cases. Hence \(\forall b_{1}\, :\, \alpha,\; b_{1} \in \{a\} \cup s \implies \neg \operatorname{Disjoint}(b,\allowbreak{} b_{1})\).
\end{proof}

\begin{theorem}[{Intersecting insert}]\label{theorem:7}
Let \(\alpha\) be a type. Fix a structure \(\mathit{inst}\) of type \(\operatorname{Semilattice\,Inf}\left(\alpha\right)\). Fix a structure \(\mathit{inst}_{1}\) of type \(\operatorname{Order\,Bot}\left(\alpha\right)\). Let \(s\) be a subset of \(\alpha\). Let \(a\) be an element of \(\alpha\). Then
\begin{center}
\(\displaystyle \operatorname{Intersecting}(\{a\} \cup s) \iff \operatorname{Intersecting}(s) \land (a \ne \operatorname{bot} \land (\forall b\, :\, \alpha,\; b \in s \implies \neg \operatorname{Disjoint}(a,\allowbreak{} b))).\)
\end{center}
\end{theorem}

\begin{proof}
Consider the following case.
\begin{proofstep}
Assume \(\operatorname{Intersecting}\left(\{a\} \cup s\right)\).
By \paperlink{premise:5365742e7375627365745f696e73657274}{Prerequisite~P\ref{premise:5365742e7375627365745f696e73657274}} with \(x = a\), we obtain \(s \subseteq \{a\} \cup s\).
By \paperlink{theorem:1}{Lemma~\ref{theorem:1}}, we obtain \(\operatorname{Intersecting}\left(s\right)\).
By \paperlink{premise:5365742e6d656d5f696e73657274}{Prerequisite~P\ref{premise:5365742e6d656d5f696e73657274}} with \(x = a\), we obtain \(a \in \{a\} \cup s\).
By \paperlink{theorem:3}{Lemma~\ref{theorem:3}}, we obtain \(a \ne \operatorname{bot}\).
For the required implication, argue as follows.
\begin{proofstep}
Fix \(x\) in \(\alpha\).
Assume \(x \in s\).
By \paperlink{premise:5365742e6d656d5f696e73657274}{Prerequisite~P\ref{premise:5365742e6d656d5f696e73657274}} with \(a\) in place of its parameter \(x\), we obtain \(a \in \{a\} \cup s\).
Use the hypothesis \(x \in s\).
By \paperlink{premise:5365742e6d656d5f696e736572745f6f665f6d656d}{Prerequisite~P\ref{premise:5365742e6d656d5f696e736572745f6f665f6d656d}} with \(y = a\), we obtain \(x \in \{a\} \cup s\).
Applying the hypothesis \(\operatorname{Intersecting}\left(\{a\} \cup s\right)\), we obtain \(\neg \operatorname{Disjoint}\left(a, x\right)\).
\end{proofstep}
Combining the two conclusions, we obtain \(a \ne \operatorname{bot} \land (\forall b\, :\, \alpha,\; b \in s \implies \neg \operatorname{Disjoint}(a,\allowbreak{} b))\).
Combining the two conclusions, we obtain \(\operatorname{Intersecting}(s) \land (a \ne \operatorname{bot} \land (\forall b\, :\, \alpha,\; b \in s \implies \neg \operatorname{Disjoint}(a,\allowbreak{} b)))\).
\end{proofstep}
Consider the following case.
\begin{proofstep}
Assume \(\operatorname{Intersecting}(s) \land (a \ne \operatorname{bot} \land (\forall b\, :\, \alpha,\; b \in s \implies \neg \operatorname{Disjoint}(a,\allowbreak{} b)))\).
Use the hypothesis \(\operatorname{Intersecting}(s) \land (a \ne \operatorname{bot} \land (\forall b\, :\, \alpha,\; b \in s \implies \neg \operatorname{Disjoint}(a,\allowbreak{} b)))\).
Taking the first conjunct, we obtain \(\operatorname{Intersecting}\left(s\right)\).
Use the hypothesis \(\operatorname{Intersecting}(s) \land (a \ne \operatorname{bot} \land (\forall b\, :\, \alpha,\; b \in s \implies \neg \operatorname{Disjoint}(a,\allowbreak{} b)))\).
Taking the second conjunct, we obtain \(a \ne \operatorname{bot} \land (\forall b\, :\, \alpha,\; b \in s \implies \neg \operatorname{Disjoint}(a,\allowbreak{} b))\).
Taking the first conjunct, we obtain \(a \ne \operatorname{bot}\).
Use the hypothesis \(\operatorname{Intersecting}(s) \land (a \ne \operatorname{bot} \land (\forall b\, :\, \alpha,\; b \in s \implies \neg \operatorname{Disjoint}(a,\allowbreak{} b)))\).
Taking the second conjunct, we obtain \(a \ne \operatorname{bot} \land (\forall b\, :\, \alpha,\; b \in s \implies \neg \operatorname{Disjoint}(a,\allowbreak{} b))\).
Taking the second conjunct, we obtain \(\forall b\, :\, \alpha,\; b \in s \implies \neg \operatorname{Disjoint}\left(a, b\right)\).
By \paperlink{theorem:6}{Lemma~\ref{theorem:6}}, we obtain \(\operatorname{Intersecting}\left(\{a\} \cup s\right)\).
\end{proofstep}
Together the two implications give the equivalence. Hence \(\operatorname{Intersecting}(\{a\} \cup s) \iff \operatorname{Intersecting}(s) \land (a \ne \operatorname{bot} \land (\forall b\, :\, \alpha,\; b \in s \implies \neg \operatorname{Disjoint}(a,\allowbreak{} b)))\).
\end{proof}

\begin{lemma}[{Intersecting if and only if pairwise not disjoint}]\label{theorem:8}
Let \(\alpha\) be a type. Fix a structure \(\mathit{inst}\) of type \(\operatorname{Semilattice\,Inf}\left(\alpha\right)\). Fix a structure \(\mathit{inst}_{1}\) of type \(\operatorname{Order\,Bot}\left(\alpha\right)\). Let \(s\) be a subset of \(\alpha\). Then
\begin{center}
\(\displaystyle \operatorname{Intersecting}(s) \iff \operatorname{Pairwise}(s,\allowbreak{} a \mapsto b \mapsto \neg \operatorname{Disjoint}(a,\allowbreak{} b)) \land s \ne \{\operatorname{bot}\}.\)
\end{center}
\end{lemma}

\begin{proof}
Consider the following case.
\begin{proofstep}
Assume \(\operatorname{Intersecting}\left(s\right)\).
For the required implication, argue as follows.
\begin{proofstep}
Fix \(a\) in \(\alpha\).
Assume \(a \in s\).
Fix \(b\) in \(\alpha\).
Assume \(b \in s\).
Assume \(a \ne b\).
Use the hypothesis \(a \in s\).
Use the hypothesis \(b \in s\).
Applying the hypothesis \(\operatorname{Intersecting}\left(s\right)\), we obtain \(\neg \operatorname{Disjoint}\left(a, b\right)\).
\end{proofstep}
For the required implication, argue as follows.
\begin{proofstep}
Assume \(s = \{\operatorname{bot}\}\).
For the required implication, argue as follows.
\begin{proofstep}
Assume \(\operatorname{Intersecting}\left(\{\operatorname{bot}\}\right)\).
By \paperlink{theorem:5}{Lemma~\ref{theorem:5}}, we obtain \(\operatorname{Intersecting}\left(\{\operatorname{bot}\}\right) \iff \operatorname{bot} \ne \operatorname{bot}\).
Use the hypothesis \(\operatorname{Intersecting}\left(\{\operatorname{bot}\}\right)\).
By definition, we obtain \(\operatorname{bot} = \operatorname{bot}\).
By the forward implication, we obtain \(\mathrm{False}\).
\end{proofstep}
Use the hypothesis \(s = \{\operatorname{bot}\}\).
Reversing the equality, we obtain \(\{\operatorname{bot}\} = s\).
Use the hypothesis \(\operatorname{Intersecting}\left(s\right)\).
Substituting equals, we obtain \(\mathrm{False}\).
\end{proofstep}
Combining the two conclusions, we obtain \(\operatorname{Pairwise}(s,\allowbreak{} a \mapsto b \mapsto \neg \operatorname{Disjoint}(a,\allowbreak{} b)) \land s \ne \{\operatorname{bot}\}\).
\end{proofstep}
Consider the following case.
\begin{proofstep}
Assume \(\operatorname{Pairwise}(s,\allowbreak{} a \mapsto b \mapsto \neg \operatorname{Disjoint}(a,\allowbreak{} b)) \land s \ne \{\operatorname{bot}\}\).
Fix \(a\) in \(\alpha\).
Assume \(a \in s\).
Fix \(b\) in \(\alpha\).
Assume \(b \in s\).
Assume \(\operatorname{Disjoint}\left(a, b\right)\).
Use the hypothesis \(\operatorname{Pairwise}(s,\allowbreak{} a \mapsto b \mapsto \neg \operatorname{Disjoint}(a,\allowbreak{} b)) \land s \ne \{\operatorname{bot}\}\).
Taking the first conjunct, we obtain \(\operatorname{Pairwise}(s,\allowbreak{} a_{1} \mapsto b_{1} \mapsto \neg \operatorname{Disjoint}(a_{1},\allowbreak{} b_{1}))\).
Use the hypothesis \(a \in s\).
Use the hypothesis \(b \in s\).
By \paperlink{premise:436c6173736963616c2e6e6f745f6e6f74}{Prerequisite~P\ref{premise:436c6173736963616c2e6e6f745f6e6f74}}, we obtain \(\neg \left(\neg \operatorname{Disjoint}\left(a, b\right)\right) \iff \operatorname{Disjoint}\left(a, b\right)\).
Use the hypothesis \(\operatorname{Disjoint}\left(a, b\right)\).
By the reverse implication, we obtain \(\neg \left(\neg \operatorname{Disjoint}\left(a, b\right)\right)\).
By \paperlink{premise:5365742e50616972776973652e6571}{Prerequisite~P\ref{premise:5365742e50616972776973652e6571}}, we obtain \(a = b\).
Use the hypothesis \(\operatorname{Pairwise}(s,\allowbreak{} a \mapsto b \mapsto \neg \operatorname{Disjoint}(a,\allowbreak{} b)) \land s \ne \{\operatorname{bot}\}\).
By \paperlink{premise:5365742e65715f73696e676c65746f6e5f6966665f756e697175655f6d656d}{Prerequisite~P\ref{premise:5365742e65715f73696e676c65746f6e5f6966665f756e697175655f6d656d}}, we obtain \(s = \{\operatorname{bot}\} \iff \operatorname{bot} \in s \land (\forall x\, :\, \alpha,\; x \in s \implies x = \operatorname{bot})\).
Use the hypothesis \(b \in s\).
Rewriting with propositional extensionality, \paperlink{premise:6469736a6f696e745f73656c66}{Prerequisite~P\ref{premise:6469736a6f696e745f73656c66}}, \(a = b\), \(\operatorname{Disjoint}\left(a, b\right)\) and substitution of equals in the preceding result gives \(\operatorname{bot} \in s\).
For the required implication, argue as follows.
\begin{proofstep}
Fix \(c\) in \(\alpha\).
Assume \(c \in s\).
By \paperlink{premise:6e6f745f6e655f696666}{Prerequisite~P\ref{premise:6e6f745f6e655f696666}}, we obtain \(\neg \left(c \ne \operatorname{bot}\right) \iff c = \operatorname{bot}\).
For the required implication, argue as follows.
\begin{proofstep}
Assume \(c \ne \operatorname{bot}\).
Use the hypothesis \(\operatorname{Pairwise}(s,\allowbreak{} a \mapsto b \mapsto \neg \operatorname{Disjoint}(a,\allowbreak{} b)) \land s \ne \{\operatorname{bot}\}\).
Use the hypothesis \(b \in s\).
Rewriting with propositional extensionality, \paperlink{premise:6469736a6f696e745f73656c66}{Prerequisite~P\ref{premise:6469736a6f696e745f73656c66}}, \(a = b\), \(\operatorname{Disjoint}\left(a, b\right)\) and substitution of equals in the preceding result gives \(\operatorname{bot} \in s\).
Use the hypothesis \(c \in s\).
Use the hypothesis \(c \ne \operatorname{bot}\).
By \paperlink{premise:4e652e73796d6d}{Prerequisite~P\ref{premise:4e652e73796d6d}}, we obtain \(\operatorname{bot} \ne c\).
By \paperlink{premise:6469736a6f696e745f626f745f6c656674}{Prerequisite~P\ref{premise:6469736a6f696e745f626f745f6c656674}}, we obtain \(\operatorname{Disjoint}\left(\operatorname{bot}, c\right)\).
Taking the first conjunct, we obtain \(\mathrm{False}\).
\end{proofstep}
By the forward implication, we obtain \(c = \operatorname{bot}\).
\end{proofstep}
Combining the two conclusions, we obtain \(\operatorname{bot} \in s \land \left(\forall x\, :\, \alpha,\; x \in s \implies x = \operatorname{bot}\right)\).
By the reverse implication, we obtain \(s = \{\operatorname{bot}\}\).
Taking the second conjunct, we obtain \(\mathrm{False}\).
\end{proofstep}
Together the two implications give the equivalence. Hence \(\operatorname{Intersecting}(s) \iff \operatorname{Pairwise}(s,\allowbreak{} a \mapsto b \mapsto \neg \operatorname{Disjoint}(a,\allowbreak{} b)) \land s \ne \{\operatorname{bot}\}\).
\end{proof}

\begin{lemma}[{Intersecting}]\label{theorem:9}
Let \(\alpha\) be a type. Fix a structure \(\mathit{inst}\) of type \(\operatorname{Semilattice\,Inf}\left(\alpha\right)\). Fix a structure \(\mathit{inst}_{1}\) of type \(\operatorname{Order\,Bot}\left(\alpha\right)\). Let \(s\) be a subset of \(\alpha\). Assume \(\operatorname{Set\mathord{:}Subsingleton}\left(s\right)\). Then \(\operatorname{Intersecting}\left(s\right) \iff s \ne \{\operatorname{bot}\}\).
\end{lemma}

\begin{proof}
By \paperlink{theorem:8}{Lemma~\ref{theorem:8}}, we obtain \(\operatorname{Intersecting}(s) \iff \operatorname{Pairwise}(s,\allowbreak{} a \mapsto b \mapsto \neg \operatorname{Disjoint}(a,\allowbreak{} b)) \land s \ne \{\operatorname{bot}\}\).
Use the hypothesis \(\operatorname{Set\mathord{:}Subsingleton}\left(s\right)\).
By \paperlink{premise:5365742e53756273696e676c65746f6e2e7061697277697365}{Prerequisite~P\ref{premise:5365742e53756273696e676c65746f6e2e7061697277697365}} with \(r = a \mapsto b \mapsto \neg \operatorname{Disjoint}\left(a, b\right)\), we obtain \(\operatorname{Pairwise}\left(s, a \mapsto b \mapsto \neg \operatorname{Disjoint}\left(a, b\right)\right)\).
By \paperlink{premise:616e645f6966665f7269676874}{Prerequisite~P\ref{premise:616e645f6966665f7269676874}}, we obtain \(\operatorname{Pairwise}(s,\allowbreak{} a \mapsto b \mapsto \neg \operatorname{Disjoint}(a,\allowbreak{} b)) \land s \ne \{\operatorname{bot}\} \iff s \ne \{\operatorname{bot}\}\).
Chaining the equivalences, we obtain \(\operatorname{Intersecting}\left(s\right) \iff s \ne \{\operatorname{bot}\}\).
\end{proof}

\begin{theorem}[{Intersecting if and only if equals empty}]\label{theorem:10}
Let \(\alpha\) be a type. Fix a structure \(\mathit{inst}\) of type \(\operatorname{Semilattice\,Inf}\left(\alpha\right)\). Fix a structure \(\mathit{inst}_{1}\) of type \(\operatorname{Order\,Bot}\left(\alpha\right)\). Assume \(\alpha\text{ has at most one element}\). Let \(s\) be a subset of \(\alpha\). Then \(\operatorname{Intersecting}\left(s\right) \iff s = \varnothing\).
\end{theorem}

\begin{proof}
Use the hypothesis \(\alpha\text{ has at most one element}\).
By \paperlink{premise:5365742e73756273696e676c65746f6e5f6f665f73756273696e676c65746f6e}{Prerequisite~P\ref{premise:5365742e73756273696e676c65746f6e5f6f665f73756273696e676c65746f6e}}, we obtain \(\operatorname{Set\mathord{:}Subsingleton}\left(s\right)\).
By \paperlink{theorem:9}{Lemma~\ref{theorem:9}}, we obtain \(\operatorname{Intersecting}\left(s\right) \iff s \ne \{\operatorname{bot}\}\).
By \paperlink{premise:6e6f745f696d705f636f6d6d}{Prerequisite~P\ref{premise:6e6f745f696d705f636f6d6d}}, we obtain \(\neg (s = \{\operatorname{bot}\}) \implies s = \varnothing \iff \neg (s = \varnothing) \implies s = \{\operatorname{bot}\}\).
For the required implication, argue as follows.
\begin{proofstep}
Assume \(\neg \left(s = \varnothing\right)\).
Use the hypothesis \(\alpha\text{ has at most one element}\).
By \paperlink{premise:5365742e73756273696e676c65746f6e5f6f665f73756273696e676c65746f6e}{Prerequisite~P\ref{premise:5365742e73756273696e676c65746f6e5f6f665f73756273696e676c65746f6e}}, we obtain \(\operatorname{Set\mathord{:}Subsingleton}\left(s\right)\).
By \paperlink{premise:5365742e6e6f6e656d7074795f6966665f6e655f656d707479}{Prerequisite~P\ref{premise:5365742e6e6f6e656d7074795f6966665f6e655f656d707479}}, we obtain \(s\ne\varnothing \iff s \ne \varnothing\).
Use the hypothesis \(\neg \left(s = \varnothing\right)\).
By the reverse implication, we obtain \(s\ne\varnothing\).
Consider the following case.
\begin{proofstep}
Fix \(a\) in \(\alpha\).
Assume \(a \in s\).
Use the hypothesis \(a \in s\).
Rewriting with \paperlink{premise:53756273696e676c65746f6e2e656c696d}{Prerequisite~P\ref{premise:53756273696e676c65746f6e2e656c696d}} (using \(\alpha\text{ has at most one element}\)) in the preceding result gives \(\operatorname{bot} \in s\).
\end{proofstep}
Using the existential hypothesis, we obtain \(\operatorname{bot} \in s\).
By \paperlink{premise:5365742e53756273696e676c65746f6e2e65715f73696e676c65746f6e5f6f665f6d656d}{Prerequisite~P\ref{premise:5365742e53756273696e676c65746f6e2e65715f73696e676c65746f6e5f6f665f6d656d}}, we obtain \(s = \{\operatorname{bot}\}\).
\end{proofstep}
By the reverse implication, we obtain \(\neg \left(s = \{\operatorname{bot}\}\right) \implies s = \varnothing\).
Consider the following case.
\begin{proofstep}
Assume \(s = \varnothing\).
By \paperlink{premise:5365742e73696e676c65746f6e5f6e6f6e656d707479}{Prerequisite~P\ref{premise:5365742e73696e676c65746f6e5f6e6f6e656d707479}} with \(\operatorname{bot}\) in place of its parameter \(a\), we obtain \(\{\operatorname{bot}\}\ne\varnothing\).
By \paperlink{premise:5365742e4e6f6e656d7074792e6e655f656d707479}{Prerequisite~P\ref{premise:5365742e4e6f6e656d7074792e6e655f656d707479}}, we obtain \(\{\operatorname{bot}\} \ne \varnothing\).
By \paperlink{premise:4e652e73796d6d}{Prerequisite~P\ref{premise:4e652e73796d6d}} with \(\mathcal{P}(\alpha)\) as its type parameter \(\alpha\), we obtain \(\varnothing \ne \{\operatorname{bot}\}\).
Use the hypothesis \(s = \varnothing\).
Reversing the equality, we obtain \(\varnothing = s\).
Substituting equals, we obtain \(s \ne \{\operatorname{bot}\}\).
\end{proofstep}
Together the two implications give the equivalence. Hence \(s \ne \{\operatorname{bot}\} \iff s = \varnothing\).
Chaining the equivalences, we obtain \(\operatorname{Intersecting}\left(s\right) \iff s = \varnothing\).
\end{proof}

\begin{theorem}[{Is upper set}]\label{theorem:11}
Let \(\alpha\) be a type. Fix a structure \(\mathit{inst}\) of type \(\operatorname{Semilattice\,Inf}\left(\alpha\right)\). Fix a structure \(\mathit{inst}_{1}\) of type \(\operatorname{Order\,Bot}\left(\alpha\right)\). Let \(s\) be a subset of \(\alpha\). Assume \(\operatorname{Intersecting}\left(s\right)\). Assume \(\forall t\, :\, \mathcal{P}(\alpha),\; \operatorname{Intersecting}(t) \implies s \subseteq t \implies s = t\). Then \(\operatorname{Is\,Upper\,Set}\left(s\right)\).
\end{theorem}

\begin{proof}
Fix \(a\) in \(\alpha\).
Fix \(b\) in \(\alpha\).
Assume \(a \le b\).
Assume \(a \in s\).
The substitution uses the following argument.
\begin{proofstep}
Fix \(c\) in \(\alpha\).
Assume \(c \in s\).
Assume \(\operatorname{Disjoint}\left(b, c\right)\).
Use the hypothesis \(a \in s\).
Use the hypothesis \(c \in s\).
Use the hypothesis \(a \le b\).
Use the hypothesis \(\operatorname{Disjoint}\left(b, c\right)\).
By \paperlink{premise:4469736a6f696e742e6d6f6e6f5f6c656674}{Prerequisite~P\ref{premise:4469736a6f696e742e6d6f6e6f5f6c656674}}, we obtain \(\operatorname{Disjoint}\left(a, c\right)\).
Applying the hypothesis \(\operatorname{Intersecting}\left(s\right)\), we obtain \(\mathrm{False}\).
\end{proofstep}
By \paperlink{premise:5365742e6d656d5f696e73657274}{Prerequisite~P\ref{premise:5365742e6d656d5f696e73657274}} with \(x = b\), we obtain \(b \in \{b\} \cup s\).
Rewriting with \(\forall t\, :\, \mathcal{P}(\alpha),\; \operatorname{Intersecting}(t) \implies s \subseteq t \implies s = t\), \paperlink{theorem:6}{Lemma~\ref{theorem:6}} (using \(\operatorname{Intersecting}\left(s\right)\)), \paperlink{premise:6d74}{Prerequisite~P\ref{premise:6d74}}, \paperlink{premise:65715f626f745f6d6f6e6f}{Prerequisite~P\ref{premise:65715f626f745f6d6f6e6f}} (using \(a \le b\)), \paperlink{theorem:3}{Lemma~\ref{theorem:3}} (using \(\operatorname{Intersecting}\left(s\right)\) and \(a \in s\)), \paperlink{premise:5365742e7375627365745f696e73657274}{Prerequisite~P\ref{premise:5365742e7375627365745f696e73657274}} and the auxiliary result established above in the preceding result gives \(b \in s\).
\end{proof}

\begin{lemma}[{Is upper set'}]\label{theorem:12}
Let \(\alpha\) be a type. Fix a structure \(\mathit{inst}\) of type \(\operatorname{Semilattice\,Inf}\left(\alpha\right)\). Fix a structure \(\mathit{inst}_{1}\) of type \(\operatorname{Order\,Bot}\left(\alpha\right)\). Let \(s\) be a finite subset of \(\alpha\). Assume \(\operatorname{Intersecting}\left(\operatorname{underlying\,set}(s)\right)\). Assume \(\forall t\, :\, \mathcal{P}_{\mathrm{fin}}(\alpha),\; \operatorname{Intersecting}(\operatorname{underlying\,set}(t)) \implies s \subseteq t \implies s = t\). Then \(\operatorname{Is\,Upper\,Set}\left(\operatorname{underlying\,set}(s)\right)\).
\end{lemma}

\begin{proof}
Fix \(a\) in \(\alpha\).
Fix \(b\) in \(\alpha\).
Assume \(a \le b\).
Assume \(a \in \operatorname{underlying\,set}(s)\).
By \paperlink{premise:46696e7365742e6d656d5f696e736572745f73656c66}{Prerequisite~P\ref{premise:46696e7365742e6d656d5f696e736572745f73656c66}} with \(b\) in place of its parameter \(a\), we obtain \(b \in \operatorname{underlying\,set}(\{b\} \cup s)\).
Rewriting with \(\forall t\, :\, \mathcal{P}_{\mathrm{fin}}(\alpha),\; \operatorname{Intersecting}(\operatorname{underlying\,set}(t)) \implies s \subseteq t \implies s = t\), \paperlink{premise:46696e7365742e636f655f696e73657274}{Prerequisite~P\ref{premise:46696e7365742e636f655f696e73657274}}, \paperlink{theorem:6}{Lemma~\ref{theorem:6}}, \(\operatorname{Intersecting}\left(\operatorname{underlying\,set}(s)\right)\), \paperlink{premise:6d74}{Prerequisite~P\ref{premise:6d74}}, \paperlink{premise:65715f626f745f6d6f6e6f}{Prerequisite~P\ref{premise:65715f626f745f6d6f6e6f}}, \(a \le b\), \paperlink{theorem:3}{Lemma~\ref{theorem:3}}, \(a \in \operatorname{underlying\,set}(s)\), \paperlink{premise:4469736a6f696e742e6d6f6e6f5f6c656674}{Prerequisite~P\ref{premise:4469736a6f696e742e6d6f6e6f5f6c656674}}, substitution of equals and \paperlink{premise:46696e7365742e7375627365745f696e73657274}{Prerequisite~P\ref{premise:46696e7365742e7375627365745f696e73657274}} in the preceding result gives \(b \in \operatorname{underlying\,set}(s)\).
\end{proof}

\section{Sets}
\begin{theorem}[{Exists membership set}]\label{theorem:13}
Let \(\alpha\) be a type. Let \(\mathcal{A}\) be a subset of \(\mathcal{P}(\alpha)\). Assume \(\operatorname{Intersecting}\left(\mathcal{A}\right)\). Let \(s\) and \(t\) be subsets of \(\alpha\). Assume \(s \in \mathcal{A}\). Assume \(t \in \mathcal{A}\). Then \(\exists a\, :\, \alpha,\; a \in s \land a \in t\).
\end{theorem}

\begin{proof}
By \paperlink{premise:5365742e6e6f745f6469736a6f696e745f696666}{Prerequisite~P\ref{premise:5365742e6e6f745f6469736a6f696e745f696666}}, we obtain \(\neg \operatorname{Disjoint}(s,\allowbreak{} t) \iff (\exists x\, :\, \alpha,\; x \in s \land x \in t)\).
Use the hypothesis \(s \in \mathcal{A}\).
Use the hypothesis \(t \in \mathcal{A}\).
Applying the hypothesis \(\operatorname{Intersecting}\left(\mathcal{A}\right)\), we obtain \(\neg \operatorname{Disjoint}\left(s, t\right)\).
By the forward implication, we obtain \(\exists a\, :\, \alpha,\; a \in s \land a \in t\).
\end{proof}

\begin{theorem}[{Exists membership finset}]\label{theorem:14}
Let \(\alpha\) be a type. Assume equality on \(\alpha\) is decidable. Let \(\mathcal{A}\) be a subset of \(\mathcal{P}_{\mathrm{fin}}(\alpha)\). Assume \(\operatorname{Intersecting}\left(\mathcal{A}\right)\). Let \(s\) and \(t\) be finite subsets of \(\alpha\). Assume \(s \in \mathcal{A}\). Assume \(t \in \mathcal{A}\). Then \(\exists a\, :\, \alpha,\; a \in s \land a \in t\).
\end{theorem}

\begin{proof}
By \paperlink{premise:5365742e6e6f745f6469736a6f696e745f696666}{Prerequisite~P\ref{premise:5365742e6e6f745f6469736a6f696e745f696666}}, we obtain \(\neg \operatorname{Disjoint}(y \mapsto y \in \operatorname{underlying\,multiset}(s),\allowbreak{} y \mapsto y \in \operatorname{underlying\,multiset}(t)) \iff (\exists x\, :\, \alpha,\; x \in (y \mapsto y \in \operatorname{underlying\,multiset}(s)) \land x \in (y \mapsto y \in \operatorname{underlying\,multiset}(t)))\).
By \paperlink{premise:46696e7365742e6469736a6f696e745f636f65}{Prerequisite~P\ref{premise:46696e7365742e6469736a6f696e745f636f65}}, we obtain \(\operatorname{Disjoint}(\operatorname{underlying\,set}(s),\allowbreak{} \operatorname{underlying\,set}(t)) \iff \operatorname{Disjoint}(s,\allowbreak{} t)\).
By \paperlink{premise:4966662e6e6f74}{Prerequisite~P\ref{premise:4966662e6e6f74}}, we obtain \(\neg \operatorname{Disjoint}(\operatorname{underlying\,set}(s),\allowbreak{} \operatorname{underlying\,set}(t)) \iff \neg \operatorname{Disjoint}(s,\allowbreak{} t)\).
Use the hypothesis \(s \in \mathcal{A}\).
Use the hypothesis \(t \in \mathcal{A}\).
Applying the hypothesis \(\operatorname{Intersecting}\left(\mathcal{A}\right)\), we obtain \(\neg \operatorname{Disjoint}\left(s, t\right)\).
By the reverse implication, we obtain \(\neg \operatorname{Disjoint}\left(\operatorname{underlying\,set}(s), \operatorname{underlying\,set}(t)\right)\).
By the forward implication, we obtain \(\exists a\, :\, \alpha,\; a \in s \land a \in t\).
\end{proof}

\begin{lemma}[{Not compl membership}]\label{theorem:15}
Let \(\alpha\) be a type. Fix a structure \(\mathit{inst}\) of type \(\operatorname{Boolean\,Algebra}\left(\alpha\right)\). Let \(s\) be a subset of \(\alpha\). Assume \(\operatorname{Intersecting}\left(s\right)\). Let \(a\) be an element of \(\alpha\). Assume \(a \in s\). Then \(\operatorname{compl}\left(a\right) \notin s\).
\end{lemma}

\begin{proof}
Assume \(\operatorname{compl}\left(a\right) \in s\).
Use the hypothesis \(a \in s\).
Use the hypothesis \(\operatorname{compl}\left(a\right) \in s\).
By \paperlink{premise:6469736a6f696e745f636f6d706c5f7269676874}{Prerequisite~P\ref{premise:6469736a6f696e745f636f6d706c5f7269676874}}, we obtain \(\operatorname{Disjoint}\left(a, \operatorname{compl}\left(a\right)\right)\).
Applying the hypothesis \(\operatorname{Intersecting}\left(s\right)\), we obtain \(\mathrm{False}\).
\end{proof}

\begin{theorem}[{Not membership}]\label{theorem:16}
Let \(\alpha\) be a type. Fix a structure \(\mathit{inst}\) of type \(\operatorname{Boolean\,Algebra}\left(\alpha\right)\). Let \(s\) be a subset of \(\alpha\). Assume \(\operatorname{Intersecting}\left(s\right)\). Let \(a\) be an element of \(\alpha\). Assume \(\operatorname{compl}\left(a\right) \in s\). Then \(a \notin s\).
\end{theorem}

\begin{proof}
Assume \(a \in s\).
Use the hypothesis \(\operatorname{compl}\left(a\right) \in s\).
Use the hypothesis \(a \in s\).
By \paperlink{premise:6469736a6f696e745f636f6d706c5f6c656674}{Prerequisite~P\ref{premise:6469736a6f696e745f636f6d706c5f6c656674}}, we obtain \(\operatorname{Disjoint}\left(\operatorname{compl}\left(a\right), a\right)\).
Applying the hypothesis \(\operatorname{Intersecting}\left(s\right)\), we obtain \(\mathrm{False}\).
\end{proof}

\begin{lemma}[{Disjoint map compl}]\label{theorem:17}
Let \(\alpha\) be a type. Fix a structure \(\mathit{inst}\) of type \(\operatorname{Boolean\,Algebra}\left(\alpha\right)\). Let \(s\) be a finite subset of \(\alpha\). Assume \(\operatorname{Intersecting}\left(\operatorname{underlying\,set}(s)\right)\). Then
\begin{center}
\(\displaystyle \operatorname{Disjoint}(s,\allowbreak{} \operatorname{map}(\operatorname{mk}(\operatorname{compl}),\allowbreak{} s)).\)
\end{center}
\end{lemma}

\begin{proof}
Fix \(a\) in \(\alpha\).
Assume \(a \in s\).
Assume \(a \in \operatorname{map}\left(\operatorname{mk}\left(\operatorname{compl}\right), s\right)\).
By \paperlink{premise:46696e7365742e6d656d5f6d6170}{Prerequisite~P\ref{premise:46696e7365742e6d656d5f6d6170}}, we obtain \(a \in \operatorname{map}(\operatorname{mk}(\operatorname{compl}),\allowbreak{} s) \iff (\exists a_{2}\, :\, \alpha,\; a_{2} \in s \land \operatorname{coe}(\operatorname{mk}(\operatorname{compl}),\allowbreak{} a_{2}) = a)\).
Use the hypothesis \(a \in \operatorname{map}\left(\operatorname{mk}\left(\operatorname{compl}\right), s\right)\).
By the forward implication, we obtain \(\exists a_{2}\, :\, \alpha,\; a_{2} \in s \land \operatorname{coe}(\operatorname{mk}(\operatorname{compl}),\allowbreak{} a_{2}) = a\).
Consider the following case.
\begin{proofstep}
Fix \(x\) in \(\alpha\).
Assume \(x \in s \land \operatorname{coe}\left(\operatorname{mk}\left(\operatorname{compl}\right), x\right) = a\).
Use the hypothesis \(x \in s \land \operatorname{coe}\left(\operatorname{mk}\left(\operatorname{compl}\right), x\right) = a\).
Consider the following case.
\begin{proofstep}
Assume \(x \in s\).
Assume \(\operatorname{coe}\left(\operatorname{mk}\left(\operatorname{compl}\right), x\right) = a\).
For the required implication, argue as follows.
\begin{proofstep}
Assume \(\operatorname{coe}\left(\operatorname{mk}\left(\operatorname{compl}\right), x\right) \in s\).
Assume \(\operatorname{coe}(\operatorname{mk}(\operatorname{compl}),\allowbreak{} x) \in \operatorname{map}(\operatorname{mk}(\operatorname{compl}),\allowbreak{} s)\).
By \paperlink{theorem:15}{Lemma~\ref{theorem:15}}, we obtain \(\mathrm{False}\).
\end{proofstep}
Use the hypothesis \(\operatorname{coe}\left(\operatorname{mk}\left(\operatorname{compl}\right), x\right) = a\).
Use the hypothesis \(a \in s\).
Use the hypothesis \(a \in \operatorname{map}\left(\operatorname{mk}\left(\operatorname{compl}\right), s\right)\).
Substituting equals, we obtain \(\mathrm{False}\).
\end{proofstep}
Using the two conjuncts, we obtain \(\mathrm{False}\).
\end{proofstep}
Using the existential hypothesis, we obtain \(\mathrm{False}\).
Rewriting with \paperlink{premise:46696e7365742e6469736a6f696e745f6c656674}{Prerequisite~P\ref{premise:46696e7365742e6469736a6f696e745f6c656674}} in the preceding result gives \(\operatorname{Disjoint}(s,\allowbreak{} \operatorname{map}(\operatorname{mk}(\operatorname{compl}),\allowbreak{} s))\).
\end{proof}

\begin{lemma}[{Cardinality at most}]\label{theorem:18}
Let \(\alpha\) be a type. Fix a structure \(\mathit{inst}\) of type \(\operatorname{Boolean\,Algebra}\left(\alpha\right)\). Fix a finite enumeration of \(\alpha\). Let \(s\) be a finite subset of \(\alpha\). Assume \(\operatorname{Intersecting}\left(\operatorname{underlying\,set}(s)\right)\). Then \(2 \cdot \left|s\right| \le \left|\alpha\right|\).
\end{lemma}

\begin{proof}
By \paperlink{premise:46696e7365742e636172645f6c655f756e6976}{Prerequisite~P\ref{premise:46696e7365742e636172645f6c655f756e6976}} with \(\operatorname{disj\,Union}(s,\allowbreak{} \operatorname{map}(\operatorname{mk}(\operatorname{compl}),\allowbreak{} s))\) in place of its parameter \(s\), we obtain \(|\operatorname{disj\,Union}(s,\allowbreak{} \operatorname{map}(\operatorname{mk}(\operatorname{compl}),\allowbreak{} s))| \le |\alpha|\).
Rewriting with \paperlink{premise:46696e7365742e636172645f6d6170}{Prerequisite~P\ref{premise:46696e7365742e636172645f6d6170}}, \paperlink{premise:46696e7365742e636172645f6469736a556e696f6e}{Prerequisite~P\ref{premise:46696e7365742e636172645f6469736a556e696f6e}}, \paperlink{theorem:17}{Lemma~\ref{theorem:17}} (using \(\operatorname{Intersecting}\left(\operatorname{underlying\,set}(s)\right)\)) and \paperlink{premise:4e61742e74776f5f6d756c}{Prerequisite~P\ref{premise:4e61742e74776f5f6d756c}} gives \(2 \cdot |s| = |\operatorname{disj\,Union}(s,\allowbreak{} \operatorname{map}(\operatorname{mk}(\operatorname{compl}),\allowbreak{} s))|\).
By \paperlink{premise:4c452e6c652e7472616e735f657127}{Prerequisite~P\ref{premise:4c452e6c652e7472616e735f657127}} with \(\mathbb{N}\) as its type parameter \(\alpha\), we obtain \(2 \cdot \left|s\right| \le \left|\alpha\right|\).
\end{proof}

\begin{lemma}[{Is max if and only if cardinality equals}]\label{theorem:19}
Let \(\alpha\) be a type. Fix a structure \(\mathit{inst}\) of type \(\operatorname{Boolean\,Algebra}\left(\alpha\right)\). Let \(\mathit{inst}_{1}\) be a proof of \(\operatorname{Nontrivial}\left(\alpha\right)\). Fix a finite enumeration of \(\alpha\). Let \(s\) be a finite subset of \(\alpha\). Assume \(\operatorname{Intersecting}\left(\operatorname{underlying\,set}(s)\right)\). Then
\begin{center}
\(\displaystyle (\forall t\, :\, \mathcal{P}_{\mathrm{fin}}(\alpha),\; \operatorname{Intersecting}(\operatorname{underlying\,set}(t)) \implies s \subseteq t \implies s = t) \iff 2 \cdot |s| = |\alpha|.\)
\end{center}
\end{lemma}

\begin{proof}
Consider the following case.
\begin{proofstep}
Assume \(\forall t\, :\, \mathcal{P}_{\mathrm{fin}}(\alpha),\; \operatorname{Intersecting}(\operatorname{underlying\,set}(t)) \implies s \subseteq t \implies s = t\).
For the required implication, argue as follows.
\begin{proofstep}
Fix \(a\) in \(\alpha\).
Consider the following case.
\begin{proofstep}
Assume \(\neg (a \in \operatorname{underlying\,set}(s) \lor \operatorname{compl}(a) \in \operatorname{underlying\,set}(s))\).
Use the hypothesis \(\neg (a \in \operatorname{underlying\,set}(s) \lor \operatorname{compl}(a) \in \operatorname{underlying\,set}(s))\).
Rewriting with \paperlink{premise:4d6174686c69622e5461637469632e507573684e65672e6e6f745f6f725f6571}{Prerequisite~P\ref{premise:4d6174686c69622e5461637469632e507573684e65672e6e6f745f6f725f6571}} in the preceding result gives \(a \notin \operatorname{underlying\,set}(s) \land \operatorname{compl}(a) \notin \operatorname{underlying\,set}(s)\).
Taking the first conjunct, we obtain \(a \notin \operatorname{underlying\,set}(s)\).
For the required implication, argue as follows.
\begin{proofstep}
Assume \(a = \operatorname{bot}\).
For the required implication, argue as follows.
\begin{proofstep}
Assume \(\operatorname{bot} \notin \operatorname{underlying\,set}(s) \land \operatorname{compl}(\operatorname{bot}) \notin \operatorname{underlying\,set}(s)\).
By \paperlink{theorem:5}{Lemma~\ref{theorem:5}}, we obtain \(\operatorname{Intersecting}\left(\{\operatorname{top}\}\right) \iff \operatorname{top} \ne \operatorname{bot}\).
Use the hypothesis \(\operatorname{Nontrivial}\left(\alpha\right)\).
By \paperlink{premise:746f705f6e655f626f74}{Prerequisite~P\ref{premise:746f705f6e655f626f74}}, we obtain \(\operatorname{top} \ne \operatorname{bot}\).
By the reverse implication, we obtain \(\operatorname{Intersecting}\left(\{\operatorname{top}\}\right)\).
Rewriting with \paperlink{premise:46696e7365742e636f655f73696e676c65746f6e}{Prerequisite~P\ref{premise:46696e7365742e636f655f73696e676c65746f6e}} in the preceding result gives \(\operatorname{Intersecting}\left(\operatorname{underlying\,set}(\{\operatorname{top}\})\right)\).
Applying the hypothesis \(\forall t\, :\, \mathcal{P}_{\mathrm{fin}}(\alpha),\; \operatorname{Intersecting}(\operatorname{underlying\,set}(t)) \implies s \subseteq t \implies s = t\), we obtain \(s \subseteq \{\operatorname{top}\} \implies s = \{\operatorname{top}\}\).
By \paperlink{premise:46696e7365742e636f655f65715f656d707479}{Prerequisite~P\ref{premise:46696e7365742e636f655f65715f656d707479}}, we obtain \(\operatorname{underlying\,set}(s) = \varnothing \iff s = \varnothing\).
By \paperlink{theorem:12}{Lemma~\ref{theorem:12}}, we obtain \(\operatorname{Is\,Upper\,Set}\left(\operatorname{underlying\,set}(s)\right)\).
By \paperlink{premise:497355707065725365742e6e6f745f746f705f6d656d}{Prerequisite~P\ref{premise:497355707065725365742e6e6f745f746f705f6d656d}}, we obtain \(\operatorname{top} \notin \operatorname{underlying\,set}(s) \iff \operatorname{underlying\,set}(s) = \varnothing\).
Use the hypothesis \(\operatorname{bot} \notin \operatorname{underlying\,set}(s) \land \operatorname{compl}(\operatorname{bot}) \notin \operatorname{underlying\,set}(s)\).
Rewriting with \paperlink{premise:636f6d706c5f626f74}{Prerequisite~P\ref{premise:636f6d706c5f626f74}} in the preceding result gives \(\operatorname{bot} \notin \operatorname{underlying\,set}(s) \land \operatorname{top} \notin \operatorname{underlying\,set}(s)\).
Taking the second conjunct, we obtain \(\operatorname{top} \notin \operatorname{underlying\,set}(s)\).
By the forward implication, we obtain \(\operatorname{underlying\,set}(s) = \varnothing\).
By the forward implication, we obtain \(s = \varnothing\).
Applying the same function to equal arguments, we obtain \((s \subseteq \{\operatorname{top}\} \implies s = \{\operatorname{top}\}) = (\varnothing \subseteq \{\operatorname{top}\} \implies \varnothing = \{\operatorname{top}\})\).
Use the previously proved fact \(s \subseteq \{\operatorname{top}\} \implies s = \{\operatorname{top}\}\).
By \paperlink{premise:46696e7365742e656d7074795f737562736574}{Prerequisite~P\ref{premise:46696e7365742e656d7074795f737562736574}} with \(\{\operatorname{top}\}\) in place of its parameter \(s\), we obtain \(\varnothing \subseteq \{\operatorname{top}\}\).
Substituting along the equality, we obtain \(\varnothing = \{\operatorname{top}\}\).
Reversing the equality, we obtain \(\{\operatorname{top}\} = \varnothing\).
By \paperlink{premise:46696e7365742e73696e676c65746f6e5f6e655f656d707479}{Prerequisite~P\ref{premise:46696e7365742e73696e676c65746f6e5f6e655f656d707479}} with \(\operatorname{top}\) in place of its parameter \(a\), we obtain \(\mathrm{False}\).
\end{proofstep}
Use the hypothesis \(a = \operatorname{bot}\).
Reversing the equality, we obtain \(\operatorname{bot} = a\).
Use the hypothesis \(\neg (a \in \operatorname{underlying\,set}(s) \lor \operatorname{compl}(a) \in \operatorname{underlying\,set}(s))\).
Rewriting with \paperlink{premise:4d6174686c69622e5461637469632e507573684e65672e6e6f745f6f725f6571}{Prerequisite~P\ref{premise:4d6174686c69622e5461637469632e507573684e65672e6e6f745f6f725f6571}} in the preceding result gives \(a \notin \operatorname{underlying\,set}(s) \land \operatorname{compl}(a) \notin \operatorname{underlying\,set}(s)\).
Substituting equals, we obtain \(\mathrm{False}\).
\end{proofstep}
For the required implication, argue as follows.
\begin{proofstep}
Fix \(b\) in \(\alpha\).
Assume \(b \in \operatorname{underlying\,set}(s)\).
Assume \(\operatorname{Disjoint}\left(a, b\right)\).
Use the hypothesis \(\neg (a \in \operatorname{underlying\,set}(s) \lor \operatorname{compl}(a) \in \operatorname{underlying\,set}(s))\).
Rewriting with \paperlink{premise:4d6174686c69622e5461637469632e507573684e65672e6e6f745f6f725f6571}{Prerequisite~P\ref{premise:4d6174686c69622e5461637469632e507573684e65672e6e6f745f6f725f6571}} in the preceding result gives \(a \notin \operatorname{underlying\,set}(s) \land \operatorname{compl}(a) \notin \operatorname{underlying\,set}(s)\).
Use the hypothesis \(\operatorname{Disjoint}\left(a, b\right)\).
By \paperlink{premise:4469736a6f696e742e6c655f636f6d706c5f6c656674}{Prerequisite~P\ref{premise:4469736a6f696e742e6c655f636f6d706c5f6c656674}}, we obtain \(b \le \operatorname{compl}\left(a\right)\).
By \paperlink{theorem:12}{Lemma~\ref{theorem:12}}, we obtain \(\operatorname{compl}\left(a\right) \in \operatorname{underlying\,set}(s)\).
Taking the second conjunct, we obtain \(\mathrm{False}\).
\end{proofstep}
By \paperlink{theorem:6}{Lemma~\ref{theorem:6}}, we obtain \(\operatorname{Intersecting}\left(\{a\} \cup \operatorname{underlying\,set}(s)\right)\).
Rewriting with \paperlink{premise:46696e7365742e636f655f696e73657274}{Prerequisite~P\ref{premise:46696e7365742e636f655f696e73657274}} in the preceding result gives \(\operatorname{Intersecting}\left(\operatorname{underlying\,set}(\{a\} \cup s)\right)\).
By \paperlink{premise:46696e7365742e7375627365745f696e73657274}{Prerequisite~P\ref{premise:46696e7365742e7375627365745f696e73657274}}, we obtain \(s \subseteq \{a\} \cup s\).
Applying the hypothesis \(\forall t\, :\, \mathcal{P}_{\mathrm{fin}}(\alpha),\; \operatorname{Intersecting}(\operatorname{underlying\,set}(t)) \implies s \subseteq t \implies s = t\), we obtain \(s = \{a\} \cup s\).
By \paperlink{premise:46696e7365742e6e655f696e736572745f6f665f6e6f745f6d656d}{Prerequisite~P\ref{premise:46696e7365742e6e655f696e736572745f6f665f6e6f745f6d656d}} with \(t = s\), we obtain \(\mathrm{False}\).
\end{proofstep}
The contrary assumption leads to a contradiction. Hence \(a \in \operatorname{underlying\,set}(s) \lor \operatorname{compl}(a) \in \operatorname{underlying\,set}(s)\).
Rewriting with \paperlink{premise:5365742e6d656d5f707265696d616765}{Prerequisite~P\ref{premise:5365742e6d656d5f707265696d616765}} in the preceding result gives \(a \in \operatorname{underlying\,set}(s) \lor a \in \operatorname{compl}^{-1}[\operatorname{underlying\,set}(s)]\).
Rewriting with \paperlink{premise:5365742e6d656d5f756e696f6e}{Prerequisite~P\ref{premise:5365742e6d656d5f756e696f6e}} in the preceding result gives \(a \in \operatorname{underlying\,set}(s) \cup \operatorname{compl}^{-1}[\operatorname{underlying\,set}(s)]\).
\end{proofstep}
By \paperlink{premise:5365742e65715f756e69765f6f665f666f72616c6c}{Prerequisite~P\ref{premise:5365742e65715f756e69765f6f665f666f72616c6c}}, we obtain \(\operatorname{underlying\,set}(s) \cup \operatorname{compl}^{-1}[\operatorname{underlying\,set}(s)] = \alpha\).
Rewriting with \paperlink{premise:5365742e696d6167655f65715f707265696d6167655f6f665f696e7665727365}{Prerequisite~P\ref{premise:5365742e696d6167655f65715f707265696d6167655f6f665f696e7665727365}} and \paperlink{premise:636f6d706c5f636f6d706c}{Prerequisite~P\ref{premise:636f6d706c5f636f6d706c}} in the preceding result gives \(\operatorname{underlying\,set}(s) \cup \operatorname{compl}[\operatorname{underlying\,set}(s)] = \alpha\).
Rewriting with \paperlink{premise:46756e6374696f6e2e456d62656464696e672e636f65466e5f6d6b}{Prerequisite~P\ref{premise:46756e6374696f6e2e456d62656464696e672e636f65466e5f6d6b}} and \paperlink{premise:636f6d706c5f696e6a656374697665}{Prerequisite~P\ref{premise:636f6d706c5f696e6a656374697665}} in the preceding result gives \(\operatorname{underlying\,set}(s) \cup \operatorname{coe}(\operatorname{mk}(\operatorname{compl}))[\operatorname{underlying\,set}(s)] = \alpha\).
Rewriting with \paperlink{premise:46696e7365742e636f655f6d6170}{Prerequisite~P\ref{premise:46696e7365742e636f655f6d6170}} in the preceding result gives \(\operatorname{underlying\,set}(s) \cup \operatorname{underlying\,set}(\operatorname{map}(\operatorname{mk}(\operatorname{compl}),\allowbreak{} s)) = \alpha\).
Rewriting with \paperlink{premise:46696e7365742e636f655f756e696f6e}{Prerequisite~P\ref{premise:46696e7365742e636f655f756e696f6e}} in the preceding result gives \(\operatorname{underlying\,set}(s \cup \operatorname{map}(\operatorname{mk}(\operatorname{compl}),\allowbreak{} s)) = \alpha\).
Rewriting with \paperlink{premise:46696e7365742e6469736a556e696f6e5f65715f756e696f6e}{Prerequisite~P\ref{premise:46696e7365742e6469736a556e696f6e5f65715f756e696f6e}} and \paperlink{theorem:17}{Lemma~\ref{theorem:17}} (using \(\operatorname{Intersecting}\left(\operatorname{underlying\,set}(s)\right)\)) in the preceding result gives \(\operatorname{underlying\,set}(\operatorname{disj\,Union}(s,\allowbreak{} \operatorname{map}(\operatorname{mk}(\operatorname{compl}),\allowbreak{} s))) = \alpha\).
Rewriting with \paperlink{premise:46696e7365742e636f655f65715f756e6976}{Prerequisite~P\ref{premise:46696e7365742e636f655f65715f756e6976}} in the preceding result gives \(\operatorname{disj\,Union}(s,\allowbreak{} \operatorname{map}(\operatorname{mk}(\operatorname{compl}),\allowbreak{} s)) = \alpha\).
Rewriting with \paperlink{premise:46696e7365742e636172645f6d6170}{Prerequisite~P\ref{premise:46696e7365742e636172645f6d6170}}, \paperlink{premise:46696e7365742e636172645f6469736a556e696f6e}{Prerequisite~P\ref{premise:46696e7365742e636172645f6469736a556e696f6e}}, \paperlink{theorem:17}{Lemma~\ref{theorem:17}} (using \(\operatorname{Intersecting}\left(\operatorname{underlying\,set}(s)\right)\)), \paperlink{premise:4e61742e74776f5f6d756c}{Prerequisite~P\ref{premise:4e61742e74776f5f6d756c}} and \(\operatorname{disj\,Union}(s,\allowbreak{} \operatorname{map}(\operatorname{mk}(\operatorname{compl}),\allowbreak{} s)) = \alpha\) gives \(2 \cdot \left|s\right| = \left|\alpha\right|\).
\end{proofstep}
Consider the following case.
\begin{proofstep}
Assume \(2 \cdot \left|s\right| = \left|\alpha\right|\).
Fix \(t\) in \(\mathcal{P}_{\mathrm{fin}}(\alpha)\).
Assume \(\operatorname{Intersecting}\left(\operatorname{underlying\,set}(t)\right)\).
Assume \(s \subseteq t\).
Use the hypothesis \(s \subseteq t\).
By \paperlink{theorem:18}{Lemma~\ref{theorem:18}}, we obtain \(2 \cdot \left|t\right| \le \left|\alpha\right|\).
Use the hypothesis \(2 \cdot \left|s\right| = \left|\alpha\right|\).
Reversing the equality, we obtain \(\left|\alpha\right| = 2 \cdot \left|s\right|\).
By \paperlink{premise:4c452e6c652e7472616e735f6571}{Prerequisite~P\ref{premise:4c452e6c652e7472616e735f6571}} with \(\mathbb{N}\) as its type parameter \(\alpha\), we obtain \(2 \cdot \left|t\right| \le 2 \cdot \left|s\right|\).
By \paperlink{premise:4e61742e74776f5f706f73}{Prerequisite~P\ref{premise:4e61742e74776f5f706f73}}, we obtain \(0 < 2\).
By \paperlink{premise:4e61742e6c655f6f665f6d756c5f6c655f6d756c5f6c656674}{Prerequisite~P\ref{premise:4e61742e6c655f6f665f6d756c5f6c655f6d756c5f6c656674}}, we obtain \(\left|t\right| \le \left|s\right|\).
By \paperlink{premise:46696e7365742e65715f6f665f7375627365745f6f665f636172645f6c65}{Prerequisite~P\ref{premise:46696e7365742e65715f6f665f7375627365745f6f665f636172645f6c65}}, we obtain \(s = t\).
\end{proofstep}
Together the two implications give the equivalence. Hence \((\forall t\, :\, \mathcal{P}_{\mathrm{fin}}(\alpha),\; \operatorname{Intersecting}(\operatorname{underlying\,set}(t)) \implies s \subseteq t \implies s = t) \iff 2 \cdot |s| = |\alpha|\).
\end{proof}

\begin{theorem}[{Exists cardinality equals}]\label{theorem:20}
Let \(\alpha\) be a type. Fix a structure \(\mathit{inst}\) of type \(\operatorname{Boolean\,Algebra}\left(\alpha\right)\). Let \(\mathit{inst}_{1}\) be a proof of \(\operatorname{Nontrivial}\left(\alpha\right)\). Fix a finite enumeration of \(\alpha\). Let \(s\) be a finite subset of \(\alpha\). Assume \(\operatorname{Intersecting}\left(\operatorname{underlying\,set}(s)\right)\). Then
\begin{center}
\(\displaystyle \exists t\, :\, \mathcal{P}_{\mathrm{fin}}(\alpha),\; s \subseteq t \land (2 \cdot |t| = |\alpha| \land \operatorname{Intersecting}(\operatorname{underlying\,set}(t))).\)
\end{center}
\end{theorem}

\begin{proof}
By \paperlink{theorem:18}{Lemma~\ref{theorem:18}}, we obtain \(2 \cdot \left|s\right| \le \left|\alpha\right|\).
Use descending induction starting at \(s\) with cardinality at most \(\left\lfloor\frac{\left|\alpha\right|}{2}\right\rfloor\). The induction hypothesis applies to each strict superset within that bound.
Consider the following case.
\begin{proofstep}
Fix \(s_{1}\) in \(\mathcal{P}_{\mathrm{fin}}(\alpha)\).
Assume as induction hypothesis that \(\forall t_{2}\, :\, \mathcal{P}_{\mathrm{fin}}(\alpha),\; |t_{2}| \le \lfloor\frac{|\alpha|}{2}\rfloor \implies s_{1} \subsetneq t_{2} \implies \operatorname{Intersecting}(\operatorname{underlying\,set}(t_{2})) \implies (\exists t\, :\, \mathcal{P}_{\mathrm{fin}}(\alpha),\; t_{2} \subseteq t \land (2 \cdot |t| = |\alpha| \land \operatorname{Intersecting}(\operatorname{underlying\,set}(t))))\).
Assume \(\left|s_{1}\right| \le \left\lfloor\frac{\left|\alpha\right|}{2}\right\rfloor\).
Assume \(\operatorname{Intersecting}\left(\operatorname{underlying\,set}(s_{1})\right)\).
Distinguish whether \(\forall t\, :\, \mathcal{P}_{\mathrm{fin}}(\alpha),\; \operatorname{Intersecting}(\operatorname{underlying\,set}(t)) \implies s_{1} \subseteq t \implies s_{1} = t\) holds.
Consider the following case.
\begin{proofstep}
Assume \(\forall t\, :\, \mathcal{P}_{\mathrm{fin}}(\alpha),\; \operatorname{Intersecting}(\operatorname{underlying\,set}(t)) \implies s_{1} \subseteq t \implies s_{1} = t\).
Choose the witness \(s_{1}\).
By \paperlink{premise:5365742e5375627365742e72666c}{Prerequisite~P\ref{premise:5365742e5375627365742e72666c}}, we obtain \((y \mapsto y \in \operatorname{underlying\,multiset}(s_{1})) \subseteq (y \mapsto y \in \operatorname{underlying\,multiset}(s_{1}))\).
By \paperlink{theorem:19}{Lemma~\ref{theorem:19}}, we obtain \((\forall t\, :\, \mathcal{P}_{\mathrm{fin}}(\alpha),\; \operatorname{Intersecting}(\operatorname{underlying\,set}(t)) \implies s_{1} \subseteq t \implies s_{1} = t) \iff 2 \cdot |s_{1}| = |\alpha|\).
Use the hypothesis \(\forall t\, :\, \mathcal{P}_{\mathrm{fin}}(\alpha),\; \operatorname{Intersecting}(\operatorname{underlying\,set}(t)) \implies s_{1} \subseteq t \implies s_{1} = t\).
By the forward implication, we obtain \(2 \cdot \left|s_{1}\right| = \left|\alpha\right|\).
Use the hypothesis \(\operatorname{Intersecting}\left(\operatorname{underlying\,set}(s_{1})\right)\).
Combining the two conclusions, we obtain \(2 \cdot |s_{1}| = |\alpha| \land \operatorname{Intersecting}(\operatorname{underlying\,set}(s_{1}))\).
Combining the two conclusions, we obtain \(s_{1} \subseteq s_{1} \land (2 \cdot |s_{1}| = |\alpha| \land \operatorname{Intersecting}(\operatorname{underlying\,set}(s_{1})))\).
This witness proves existence. Hence \(\exists t\, :\, \mathcal{P}_{\mathrm{fin}}(\alpha),\; s_{1} \subseteq t \land (2 \cdot |t| = |\alpha| \land \operatorname{Intersecting}(\operatorname{underlying\,set}(t)))\).
\end{proofstep}
Consider the following case.
\begin{proofstep}
Assume \(\neg (\forall t\, :\, \mathcal{P}_{\mathrm{fin}}(\alpha),\; \operatorname{Intersecting}(\operatorname{underlying\,set}(t)) \implies s_{1} \subseteq t \implies s_{1} = t)\).
The substitution uses the following argument.
\begin{proofstep}
Fix \(t\) in \(\mathcal{P}_{\mathrm{fin}}(\alpha)\).
By \paperlink{premise:4d6174686c69622e5461637469632e507573684e65672e6e6f745f696d706c6965735f6571}{Prerequisite~P\ref{premise:4d6174686c69622e5461637469632e507573684e65672e6e6f745f696d706c6965735f6571}} with \(p = \operatorname{Intersecting}\left(\operatorname{underlying\,set}(t)\right)\) and \(q = s_{1} \subseteq t \implies s_{1} = t\), we obtain \((\neg (\operatorname{Intersecting}(\operatorname{underlying\,set}(t)) \implies s_{1} \subseteq t \implies s_{1} = t)) = (\operatorname{Intersecting}(\operatorname{underlying\,set}(t)) \land \neg (s_{1} \subseteq t \implies s_{1} = t))\).
By \paperlink{premise:4d6174686c69622e5461637469632e507573684e65672e6e6f745f696d706c6965735f6571}{Prerequisite~P\ref{premise:4d6174686c69622e5461637469632e507573684e65672e6e6f745f696d706c6965735f6571}} with \(p = s_{1} \subseteq t\) and \(q = s_{1} = t\), we obtain \((\neg (s_{1} \subseteq t \implies s_{1} = t)) = (s_{1} \subseteq t \land \neg (s_{1} = t))\).
Applying the same function to equal arguments, we obtain \((\operatorname{Intersecting}(\operatorname{underlying\,set}(t)) \land \neg (s_{1} \subseteq t \implies s_{1} = t)) = (\operatorname{Intersecting}(\operatorname{underlying\,set}(t)) \land (s_{1} \subseteq t \land \neg (s_{1} = t)))\).
Chaining the equalities, we obtain \((\neg (\operatorname{Intersecting}(\operatorname{underlying\,set}(t)) \implies s_{1} \subseteq t \implies s_{1} = t)) = (\operatorname{Intersecting}(\operatorname{underlying\,set}(t)) \land (s_{1} \subseteq t \land \neg (s_{1} = t)))\).
\end{proofstep}
Use the hypothesis \(\neg (\forall t\, :\, \mathcal{P}_{\mathrm{fin}}(\alpha),\; \operatorname{Intersecting}(\operatorname{underlying\,set}(t)) \implies s_{1} \subseteq t \implies s_{1} = t)\).
Rewriting with \paperlink{premise:4d6174686c69622e5461637469632e507573684e65672e6e6f745f666f72616c6c5f6571}{Prerequisite~P\ref{premise:4d6174686c69622e5461637469632e507573684e65672e6e6f745f666f72616c6c5f6571}}, \paperlink{premise:66756e657874}{Prerequisite~P\ref{premise:66756e657874}} and the identity established above in the preceding result gives \(\exists t\, :\, \mathcal{P}_{\mathrm{fin}}(\alpha),\; \operatorname{Intersecting}(\operatorname{underlying\,set}(t)) \land (s_{1} \subseteq t \land s_{1} \ne t)\).
Consider the following case.
\begin{proofstep}
Fix \(t\) in \(\mathcal{P}_{\mathrm{fin}}(\alpha)\).
Assume \(\operatorname{Intersecting}(\operatorname{underlying\,set}(t)) \land (s_{1} \subseteq t \land s_{1} \ne t)\).
Use the hypothesis \(\operatorname{Intersecting}(\operatorname{underlying\,set}(t)) \land (s_{1} \subseteq t \land s_{1} \ne t)\).
Consider the following case.
\begin{proofstep}
Assume \(\operatorname{Intersecting}\left(\operatorname{underlying\,set}(t)\right)\).
Assume \(s_{1} \subseteq t \land s_{1} \ne t\).
For the required implication, argue as follows.
\begin{proofstep}
Fix \(u\) in \(\mathcal{P}_{\mathrm{fin}}(\alpha)\).
By \paperlink{premise:46696e7365742e696e737449735472616e73537562736574}{Prerequisite~P\ref{premise:46696e7365742e696e737449735472616e73537562736574}}, we obtain \(\operatorname{Is\,Trans}(\mathcal{P}_{\mathrm{fin}}(\alpha),\allowbreak{} \mathit{x1} \mapsto \mathit{x2} \mapsto \mathit{x1} \subseteq \mathit{x2})\).
Use the hypothesis \(s_{1} \subseteq t \land s_{1} \ne t\).
Taking the first conjunct, we obtain \(s_{1} \subseteq t\).
By \paperlink{premise:4861735375627365742e5375627365742e7472616e73}{Prerequisite~P\ref{premise:4861735375627365742e5375627365742e7472616e73}} with \(\mathcal{P}_{\mathrm{fin}}(\alpha)\) as its type parameter \(\alpha\), we obtain \(t \subseteq u \implies s_{1} \subseteq u\).
By \paperlink{premise:416e642e696d705f6c656674}{Prerequisite~P\ref{premise:416e642e696d705f6c656674}}, we obtain \(t \subseteq u \land (2 \cdot |u| = |\alpha| \land \operatorname{Intersecting}(\operatorname{underlying\,set}(u))) \implies s_{1} \subseteq u \land (2 \cdot |u| = |\alpha| \land \operatorname{Intersecting}(\operatorname{underlying\,set}(u)))\).
\end{proofstep}
By \paperlink{theorem:18}{Lemma~\ref{theorem:18}}, we obtain \(2 \cdot \left|t\right| \le \left|\alpha\right|\).
Rewriting with \paperlink{premise:4e61742e6d756c5f636f6d6d}{Prerequisite~P\ref{premise:4e61742e6d756c5f636f6d6d}} in the preceding result gives \(\left|t\right| \cdot 2 \le \left|\alpha\right|\).
Rewriting with \paperlink{premise:4e61742e6c655f6469765f6966665f6d756c5f6c65}{Prerequisite~P\ref{premise:4e61742e6c655f6469765f6966665f6d756c5f6c65}} and \paperlink{premise:4e61742e74776f5f706f73}{Prerequisite~P\ref{premise:4e61742e74776f5f706f73}} in the preceding result gives \(\left|t\right| \le \left\lfloor\frac{\left|\alpha\right|}{2}\right\rfloor\).
By \paperlink{premise:46696e7365742e696e737449734e6f6e7374726963745374726963744f7264657253756273657453537562736574}{Prerequisite~P\ref{premise:46696e7365742e696e737449734e6f6e7374726963745374726963744f7264657253756273657453537562736574}}, we obtain \(\operatorname{Is\,Nonstrict\,Strict\,Order}(\mathcal{P}_{\mathrm{fin}}(\alpha),\allowbreak{} \mathit{x1} \mapsto \mathit{x2} \mapsto \mathit{x1} \subseteq \mathit{x2},\allowbreak{} \mathit{x1} \mapsto \mathit{x2} \mapsto \mathit{x1} \subsetneq \mathit{x2})\).
By \paperlink{premise:46696e7365742e696e73744973416e746973796d6d537562736574}{Prerequisite~P\ref{premise:46696e7365742e696e73744973416e746973796d6d537562736574}}, we obtain \(\operatorname{Is\,Antisymm}(\mathcal{P}_{\mathrm{fin}}(\alpha),\allowbreak{} \mathit{x1} \mapsto \mathit{x2} \mapsto \mathit{x1} \subseteq \mathit{x2})\).
By \paperlink{premise:737375627365745f6966665f7375627365745f6e65}{Prerequisite~P\ref{premise:737375627365745f6966665f7375627365745f6e65}} with \(\mathcal{P}_{\mathrm{fin}}(\alpha)\) as its type parameter \(\alpha\), we obtain \(s_{1} \subsetneq t \iff s_{1} \subseteq t \land s_{1} \ne t\).
Use the hypothesis \(s_{1} \subseteq t \land s_{1} \ne t\).
By the reverse implication, we obtain \(s_{1} \subsetneq t\).
Use the hypothesis \(\operatorname{Intersecting}\left(\operatorname{underlying\,set}(t)\right)\).
Applying the induction hypothesis, we obtain \(\exists t_{1}\, :\, \mathcal{P}_{\mathrm{fin}}(\alpha),\; t \subseteq t_{1} \land (2 \cdot |t_{1}| = |\alpha| \land \operatorname{Intersecting}(\operatorname{underlying\,set}(t_{1})))\).
By \paperlink{premise:4578697374732e696d70}{Prerequisite~P\ref{premise:4578697374732e696d70}} with \(\mathcal{P}_{\mathrm{fin}}(\alpha)\) as its type parameter \(\alpha\), we obtain \(\exists a\, :\, \mathcal{P}_{\mathrm{fin}}(\alpha),\; s_{1} \subseteq a \land (2 \cdot |a| = |\alpha| \land \operatorname{Intersecting}(\operatorname{underlying\,set}(a)))\).
\end{proofstep}
Using the two conjuncts, we obtain \(\exists t_{1}\, :\, \mathcal{P}_{\mathrm{fin}}(\alpha),\; s_{1} \subseteq t_{1} \land (2 \cdot |t_{1}| = |\alpha| \land \operatorname{Intersecting}(\operatorname{underlying\,set}(t_{1})))\).
\end{proofstep}
Using the existential hypothesis, we obtain \(\exists t\, :\, \mathcal{P}_{\mathrm{fin}}(\alpha),\; s_{1} \subseteq t \land (2 \cdot |t| = |\alpha| \land \operatorname{Intersecting}(\operatorname{underlying\,set}(t)))\).
\end{proofstep}
The proposition or its negation holds. Hence \(\exists t\, :\, \mathcal{P}_{\mathrm{fin}}(\alpha),\; s_{1} \subseteq t \land (2 \cdot |t| = |\alpha| \land \operatorname{Intersecting}(\operatorname{underlying\,set}(t)))\).
\end{proofstep}
Use the previously proved fact \(2 \cdot \left|s\right| \le \left|\alpha\right|\).
Rewriting with \paperlink{premise:4e61742e6d756c5f636f6d6d}{Prerequisite~P\ref{premise:4e61742e6d756c5f636f6d6d}} in the preceding result gives \(\left|s\right| \cdot 2 \le \left|\alpha\right|\).
Rewriting with \paperlink{premise:4e61742e6c655f6469765f6966665f6d756c5f6c65}{Prerequisite~P\ref{premise:4e61742e6c655f6469765f6966665f6d756c5f6c65}} and \paperlink{premise:4e61742e74776f5f706f73}{Prerequisite~P\ref{premise:4e61742e74776f5f706f73}} in the preceding result gives \(\left|s\right| \le \left\lfloor\frac{\left|\alpha\right|}{2}\right\rfloor\).
Use the hypothesis \(\operatorname{Intersecting}\left(\operatorname{underlying\,set}(s)\right)\).
By descending induction under the cardinality bound, we obtain \(\exists t\, :\, \mathcal{P}_{\mathrm{fin}}(\alpha),\; s \subseteq t \land (2 \cdot |t| = |\alpha| \land \operatorname{Intersecting}(\operatorname{underlying\,set}(t)))\).
\end{proof}

\end{document}
